% =====================================================================
% Resilient Last-Mile Allocation in Humanitarian Supply Chains:
% A Reproducible Many-Objective Fuzzy Framework
%
% Submission target: Applied Sciences (MDPI), Special Issue
%   "Innovations in Supply Chain Resilience"
%   (Guest Editors: Agnusdei, Silvestri, Di Pietro)
%
% Tier 1 skeleton — section structure + 1-sentence placeholders per
% subsection per book-implementation-plan.md. To swap in MDPI's official
% applsci-mdpi.cls before submission:
%   1. Download from https://www.mdpi.com/authors/latex into ./Definitions/
%   2. Replace the documentclass line below with the MDPI version
%   3. Move the article metadata into MDPI's macros (\Title, \Author, …)
% =====================================================================

\documentclass[11pt,a4paper]{article}

\usepackage[utf8]{inputenc}
\usepackage[T1]{fontenc}
\usepackage[final,expansion=true,protrusion=true]{microtype}
\usepackage{newtxtext,newtxmath}
\usepackage{graphicx}
\usepackage{booktabs}
\usepackage{tabularx}
\usepackage{multirow}
\usepackage{amsmath}
\let\Bbbk\relax % newtxmath predefines \Bbbk; avoid amssymb redefinition clash
\usepackage{amssymb}
\usepackage{xcolor}
\usepackage[colorlinks=true,linkcolor=blue!50!black,citecolor=blue!50!black]{hyperref}
\usepackage{xurl}
\usepackage[margin=2.5cm]{geometry}
\usepackage{caption}
\usepackage{subcaption}
\usepackage{authblk}
\usepackage{titlesec}

\titleformat{\section}{\Large\bfseries}{\thesection.}{0.5em}{}
\titleformat{\subsection}{\large\bfseries}{\thesubsection}{0.5em}{}
\graphicspath{{../figures/}}

\title{Resilient Last-Mile Allocation in Humanitarian Supply Chains:\\
       A Reproducible Many-Objective Fuzzy Framework}

\author[1]{Peyman Rabiei}
\author[1]{Daniel Arias-Aranda}
\author[2,*]{Vladimir Stantchev}
\affil[1]{Faculty of Economics and Business, University of Granada, Spain}
\affil[2]{Institute of Information Systems, SRH Berlin University, Germany}
\affil[*]{Correspondence: \href{mailto:stantchev@computer.org}{stantchev@computer.org}}

\date{}

\begin{document}
\maketitle

% (Extended-version disclosure removed: the earlier manuscript is unpublished;
% this paper is the primary publication of the model.)

% =====================================================================
\begin{abstract}
\noindent
Disaster frequency is rising and response-phase planners face a binding
constraint: relief capacity is always smaller than need, so directing
affected people to relief centres is a triage decision under time
pressure. Published optimisation models for this problem use
operations-research vocabulary (cost, distance, equity); civil-protection
planners use supply-chain-resilience vocabulary (the 4R framework of
Bruneau et al.). We close the gap. We develop a four-objective humanitarian
last-mile allocation model whose objectives map one-to-one onto the 4R
components --- Resourcefulness, Robustness, Rapidity, Redundancy ---
evaluated by Mamdani fuzzy inference systems and solved by NSGA-II, NRGA,
and NSGA-III. On 540 runs across three problem sizes the four-objective
formulation surfaces independent Robustness and Rapidity axes that the
baseline three-objective formulation fuses into one --- a visibility gain
for the decision-maker, established by rank-correlation analysis.
NSGA-III is 32\% faster than NSGA-II at every size with comparable
Pareto-front quality, the HV ranking depending on the reference-point scheme. Weight-perturbation sweeps (100
Latin-Hypercube samples, $\pm 20\%$) confirm the recommendations are
robust under expert-elicitation noise at every problem size; rule-base
deletion sweeps confirm the same for the three larger fuzzy systems,
with the three-rule Rapidity classifier flagged as a structural
sensitivity hot-spot. The implementation is released as an
MIT-licensed Python reference at
\url{https://github.com/presidio-v/presidio-hardened-vol-assign}.
\end{abstract}

\noindent\textbf{Keywords:}
supply chain resilience; humanitarian logistics; multi-objective optimization;
fuzzy inference systems; NSGA-III; reproducible research; 4R framework

% =====================================================================
\section{Introduction}\label{sec:intro}
% TIER 1 PLACEHOLDER — target ~1{,}200 words.
% Lead with disaster-frequency claim (CRED 2023). Pivot fast to humanitarian
% supply chain resilience. Bruneau's 4R named in second/third paragraph.
% Bounded problem statement (lift from explore/problem-statement.md).
% Five contributions, listed numbered.
% Closing paragraph: paper structure pointer.

The 2023 EM-DAT record lists 399 natural disasters worldwide, claiming
86{,}473 lives and inflicting US\$202.7\,bn in economic damage
\cite{cred2023, wirtz2014}. The arithmetic exceeds the 20-year baseline of
64{,}148 deaths and is projected to keep rising. Behind the aggregate
numbers sits a more compressed timeline: the response-phase mortality
literature places the bulk of preventable deaths inside the first 12 hours
after onset \cite{avishan2023}, the window in which civil-protection
planners must direct affected people to relief centres before the most
vulnerable run out of resources or infrastructure damage cuts the routes
that remain.

The decision is not whether to direct people; it is whom to direct, where
to send them, and what to accept giving up. Disaster relief in the first
hours always serves fewer than need help \cite{krishna2021, behl2019}.
Three structural concerns shape every direction decision: prioritising the
most vulnerable, choosing routes that can actually deliver despite
infrastructure damage, and balancing the load across centres so that no
single relief site exhausts before the rest. None of the three reduces to
the others; the practitioner balances all three under time pressure with
incomplete information.

Humanitarian supply-chain research has accumulated tools for this decision
class \cite{christopher2004, tukamuhabwa2015, behl2019}. Most published
allocation models, including the recent literature in this venue's
community, frame their objectives in operations-research vocabulary:
travel distance, total cost, response time, head-count of unmet demand.
The vocabulary is well-suited to optimisation but it is not the vocabulary
the policy and planning literature uses. Bruneau et al.'s 4R framework
\cite{bruneau2003} --- Robustness, Redundancy, Resourcefulness, Rapidity ---
has been the dominant resilience taxonomy in disaster risk reduction since
2003 and has spread through commercial supply-chain resilience research
\cite{christopher2004, hosseini2016}. Civil-protection planners, FEMA
documents, EU Civil Protection Mechanism plans, and academic
supply-chain-resilience reviews all use 4R vocabulary. Optimisation models
that emit results in cost/distance space force the practitioner to
translate; models that emit results in 4R space let the practitioner read
the Pareto front directly.

This paper develops a four-objective humanitarian last-mile allocation
model in which each objective maps one-to-one onto a 4R component. Three
of the four objectives build on the authors' recent FIS-MOEA line of work
\cite{rabiei2021, rabiei2023}; the fourth is new and is the
result of splitting a baseline transportation-infeasibility index
into separate Robustness and Rapidity components. We solve the model with
NSGA-II, NRGA, and NSGA-III, evaluate it under rule-base and weight
sensitivity analyses, and release an open-source Python reference
implementation that any reviewer or practitioner can run on their own
laptop in minutes. The combination addresses a gap the related-work survey
documents in Section~\ref{sec:related}: to the best of our knowledge, the
recent FIS-MOEA disaster-allocation literature does not include a model
that organises its objectives around a named resilience framework and
also ships a reproducible artifact.

\subsection*{Problem statement and contributions}
% Lift bounded problem statement from explore/problem-statement.md.

The contribution of this paper is a reframing and a tool, not a new
optimisation operator. We do not propose a new selection, crossover, or
mutation mechanism; we reorganise an existing class of allocation models
around a recognised resilience vocabulary and ship an implementation that
others can run and check. Concretely, the paper makes five contributions to
the literature on humanitarian last-mile allocation and supply chain
resilience:
\begin{enumerate}
    \item The first explicit 4R$\leftrightarrow$objective mapping in FIS-MOEA disaster
          allocation, formalised in Section~\ref{sec:resilience-mapping}.
    \item A four-objective formulation that splits the transportation
          infeasibility level into separate Robustness (TRD) and Rapidity (RPD)
          components, exposing trade-offs the three-objective formulation
          cannot represent (Section~\ref{sec:tradeoff-analysis}).
    \item An empirical comparison of three established solvers --- NSGA-II,
          NRGA, and NSGA-III --- on the four-objective humanitarian allocation
          problem, used to characterise the model rather than to advance the
          solvers (Section~\ref{sec:performance}).
    \item Rule-base and weight-perturbation sensitivity analyses establishing
          robustness of the recommendations under expert-elicitation noise
          (Section~\ref{sec:sensitivity}).
    \item An open-source Python reference implementation released with a
          citable Zenodo DOI, closing the FIS-MOEA reproducibility gap in this
          line of work (DOI:~\url{https://doi.org/10.5281/zenodo.21083544}).
\end{enumerate}

The implicit hypothesis behind the reframing is that organising the objective
set around the 4R framework buys the practitioner something measurable, not
just a relabelling. We make it explicit and decompose it into four testable
hypotheses, each resolved by the experiments in
Sections~\ref{sec:performance}--\ref{sec:sensitivity}:
\begin{description}
    \item[H1] Splitting the fused transportation index into separate Robustness
          (TRD) and Rapidity (RPD) objectives exposes trade-off information that
          the three-objective formulation cannot represent
          (Section~\ref{sec:tradeoff-analysis}).
    \item[H2] At four objectives, NSGA-III improves on NSGA-II in Pareto-front
          quality (H2a) or in computational cost (H2c), the boundary at which
          reference-point methods are expected to begin paying off
          (Section~\ref{sec:perf-algorithms}).
    \item[H3] The model's recommendations are robust to realistic
          expert-elicitation noise, both in the fuzzy rule bases (H3a) and in
          the elicited weights (H3b) (Section~\ref{sec:sensitivity}).
    \item[H4] The NSGA-II-versus-NRGA cost ranking reported by the earlier
          MATLAB implementation of this model reproduces in an independent
          Python implementation (Section~\ref{sec:perf-algorithms}).
\end{description}

The remainder is organised as follows. Section~\ref{sec:related} reviews
supply-chain-resilience theory, humanitarian last-mile allocation, FIS-MOEA
hybrids, and many-objective evolutionary algorithms. Section~\ref{sec:model}
formalises the model and introduces the 4R$\leftrightarrow$objective mapping.
Section~\ref{sec:algorithms} describes the algorithm design including the
NSGA-III reference-point construction. Section~\ref{sec:performance} reports
the algorithm comparison; Section~\ref{sec:tradeoff-analysis} the projection
analysis testing whether the 4-obj split exposes new trade-offs;
Section~\ref{sec:sensitivity} the rule-base and weight sensitivity studies.
Section~\ref{sec:discussion} discusses implications for theory and practice;
Section~\ref{sec:conclusion} concludes.

% =====================================================================
\section{Related Work and Theoretical Background}\label{sec:related}

The literature relevant to this paper sits across four bodies of work:
supply-chain-resilience theory; humanitarian last-mile allocation models;
fuzzy-inference / multi-objective evolutionary hybrids in disaster
operations; and many-objective evolutionary algorithms. We surveyed
each across Scopus, Web of Science, and IEEE Xplore using the term sets
documented per subsection below, with priority given to peer-reviewed
journal articles published between 2014 and 2025 and to landmark
foundational works regardless of date. The four sub-surveys converge on
a common gap that the present paper addresses, which we state at the end
of the section. We do not claim a systematic-review-grade exhaustive
coverage; the relative-work negative claims throughout this paper are
qualified accordingly.

\subsection{Supply chain resilience and the 4R framework}\label{sec:related-scr}

The 4R framework originates in Bruneau et al.'s seismic-resilience
quantification \cite{bruneau2003}: a community is resilient when it exhibits
\emph{Robustness} (capacity to withstand a shock without significant loss
of function), \emph{Redundancy} (substitutability of system elements),
\emph{Resourcefulness} (capacity to identify and mobilise responses), and
\emph{Rapidity} (capacity to act in time). The four components were
proposed for community infrastructure but the vocabulary has spread
\cite{hosseini2016}: it is the dominant resilience taxonomy in disaster
risk reduction, civil-protection planning documents, and the broader
SCM-resilience literature.

Christopher and Peck \cite{christopher2004} extended the resilience lens
to commercial supply chains, naming flexibility, redundancy, and
collaboration as the operational levers. Pettit, Fiksel, and Croxton
\cite{pettit2010, pettit2013} developed an assessment-tool framework
that pairs SCR capabilities against vulnerability factors. Tukamuhabwa
et al.\ \cite{tukamuhabwa2015} and Kamalahmadi and Parast
\cite{kamalahmadi2016} review more than a decade of SCR work and
consolidate it into taxonomies that map cleanly back onto Bruneau's
4R; Wieland \cite{wieland2021} reframes SCR as a continuous adaptive
process rather than a static property, but the underlying 4R
vocabulary persists. Hosseini et al.\ \cite{hosseini2016} provide the
most-cited quantitative review, distinguishing measures by what they
capture (probabilistic vs.\ deterministic) and how they aggregate
(component-wise vs.\ system-wide). Behl and Dutta \cite{behl2019}
survey humanitarian SCR specifically and identify the absence of
computational allocation models that operationalise the 4R framework
as an open research direction.

Two observations matter for the present work. First, the 4R framework is
operational, not philosophical: each component has tested measurement
protocols in the disaster-management literature, and the vocabulary is
shared between researchers and practitioners. Second, the operationalisation
of 4R \emph{inside computational allocation models} --- as opposed to
post-hoc analysis of a model's outputs --- is rare. We address that gap
directly in Section~\ref{sec:resilience-mapping}.

\subsection{Humanitarian last-mile allocation}\label{sec:related-allocation}

The classical humanitarian-allocation literature is dominated by cost,
time, and distance objectives \cite{altay2006, galindo2013, caunhye2012,
anayaarenas2014}. Tzeng et al.\ \cite{tzeng2007} present a multi-objective
relief-delivery model with cost and time; Sheu \cite{sheu2007} formulates
emergency-logistics priorities and constraints from the field-operations
perspective. Abounacer et al.\ \cite{abounacer2014} formalise
multi-objective location-transportation with cost minimisation and
uncovered-demand penalty. Gama et al.\ \cite{gama2016} introduce a
multi-period flood-shelter allocation model with evacuation-order timing
and distance objectives. Zhao et al.\ \cite{zhao2017} integrate location
and allocation with travel-distance and coverage objectives. Balcik and
Beamon \cite{balcik2008} introduce facility-location models for
humanitarian relief; Pedraza-Martinez et al.\ \cite{pedrazamartinez2011}
study fleet management in humanitarian operations. The Holgu{\'\i}n-Veras
et al.\ \cite{holguinveras2012} characterisation of post-disaster
logistics' unique features --- including the displaced-people allocation
problem central to this paper --- frames the broader research programme.
Sahebjamnia et al.\ \cite{sahebjamnia2017} build a hybrid decision
support system for humanitarian relief chains; Goldschmidt and Kumar
\cite{goldschmidt2016} review the field with operational-research
emphasis.

Equity entered the literature later. Soghrati Ghasbeh et al.
\cite{soghrati2022} develop a robust multi-stage relief-distribution model
with explicit fairness terms; Krishna et al. \cite{krishna2021} document
the practitioner-side experience of inequitable relief in the 2015 South
Indian floods, motivating the equity turn. Across this body of work the
shared characteristic is that objectives are crisp quantitative measures
(seconds, kilometres, dollars, head-counts), not qualitative indices that
encode expert reasoning under uncertainty.

Within this surveyed literature, the use of a recognised resilience
framework to organise the objective set --- particularly the 4R framework
\cite{bruneau2003} --- is rare. Models include resilience as a constraint
or a robustness term, but to the best of our knowledge none of the works
above structures its objectives as an explicit bijection with the four 4R
components. The framing remains Operations
Research's cost/time/equity vocabulary; the SCR community's vocabulary sits
unused on the next shelf over.

\subsection{FIS-MOEA hybrids in disaster operations}\label{sec:related-fis}

Mamdani's original fuzzy-controller paper \cite{mamdani1975} established
that linguistic if-then rules combined with triangular and trapezoidal
membership functions could replace tuned crisp models in domains where
expert judgement is the canonical input. The pattern fits humanitarian
operations management precisely: practitioners reason in linguistic
categories (``severely injured,'' ``partially blocked,'' ``rapidly
depleting''), and translating those to crisp scalars discards
practitioner-relevant information.

The FIS-MOEA hybrid pattern in disaster operations was introduced into
this line of work by Rabiei and Arias-Aranda \cite{rabiei2021} for
post-disaster vehicle routing and relief supply distribution, and extended
to volunteer-to-vacancy assignment in Rabiei, Arias-Aranda, and Stantchev
\cite{rabiei2023}. Both pair Mamdani FIS with NSGA-II and NRGA. The present
paper adapts the pattern to the people-to-relief-centre allocation problem,
developing the three fuzzy-aggregated indices (FLPP, TFL, CABL) and their
corresponding objective functions (ULPP, TIL, CAIL) as the baseline
three-objective formulation that the four-objective model extends. Adjacent fuzzy or stochastic disaster-allocation work
includes Bahmani et al.\ \cite{bahmani2018} on fuzzy multi-objective
routing under damaged networks and Moreno et al.\ \cite{moreno2018} on
two-stage stochastic relief-supply chains with social-equity
objectives; neither operationalises a named resilience framework.

Three observations about this line of work motivate the present extension.
First, the FIS-MOEA disaster-line papers we surveyed do not reference the
4R framework or any other named resilience taxonomy: the objective
functions are introduced on their own terms, with intuitive
justification, not as operationalisations of an existing theory. Second,
the baseline Transportation Infeasibility Level (TIL) fuses
Robustness (road condition, hazard) with Rapidity (travel duration) into
a single index --- which prevents the practitioner from articulating
Robustness/Rapidity trade-offs separately. Third, none of the FIS-MOEA
papers we surveyed releases a reproducible artifact: the prior work in
this line uses MATLAB without an accompanying open implementation. This combination
is the gap we address in Sections~\ref{sec:model}--\ref{sec:performance}.

\subsection{Many-objective evolutionary algorithms}\label{sec:related-moea}

NSGA-II \cite{deb2002} remains the most-cited multi-objective evolutionary
algorithm for two- and three-objective problems. The combination of
non-dominated sorting and crowding-distance tie-breaking is well-suited to
the regime where objectives are few enough that the Pareto front retains a
clear geometric structure and crowding distance carries usable diversity
information.

Reference-point methods entered the literature for many-objective settings
(typically $\ge 4$ objectives) where crowding distance becomes
information-poor \cite{ishibuchi2008, li2015many, coello2007}. Deb and
Jain's NSGA-III \cite{deb2014nsga3} replaces crowding distance with
niching against a structured reference set, laid down on the normalised
objective simplex via the Das-Dennis construction. Decomposition-based
methods such as MOEA/D \cite{zhang2007moead} and reference-vector
methods such as RVEA \cite{cheng2016rvea} pursue the same goal through
different mechanisms; we discuss them as candidates for future
algorithmic comparison in Section~\ref{sec:disc-future}. The
reference-point machinery scales gracefully with objective count and
is the modern default for many-objective problems. Recent runtime
analysis by Zheng and Doerr \cite{zheng2024nsga2} quantifies NSGA-II's
degradation as a function of objective count and provides the
theoretical underpinning for the empirical observation that
reference-point methods take over above about four objectives.

NRGA \cite{omar2008nrga} sits between NSGA-II and NSGA-III conceptually:
it inherits NSGA-II's non-dominated sorting but replaces the
crowding-distance tie-breaker with uniform random sampling within the
boundary front. The motivation is that crowding distance encodes a
particular notion of diversity (gap between consecutive solutions in
objective space) that may not match the practitioner's notion; random
sampling avoids that prior at the cost of dropping some non-dominated
solutions. NRGA is the second algorithm used in the FIS-MOEA disaster
line \cite{rabiei2021, rabiei2023} and we retain it for
continuity with that line.

This paper's algorithmic contribution is to add NSGA-III
to the comparison, putting it on a four-objective humanitarian allocation
problem precisely at the boundary where reference-point methods are
expected to start paying off. The empirical findings of
Section~\ref{sec:performance} show that the boundary is crossed for CPU
time but not for hypervolume --- a finding that maps onto the runtime
analysis of \cite{zheng2024nsga2} and is the substance of the H2 result.

\subsection{Gap}\label{sec:related-gap}

The combination of (i) fuzzy-aggregated objectives (quantitative inputs,
linguistic outputs) organised around a named resilience framework, (ii) a
four-objective formulation that
splits Robustness and Rapidity into separate axes following the 4R
taxonomy, (iii) NSGA-III alongside NSGA-II and NRGA on a humanitarian
last-mile allocation problem, (iv) systematic rule-base and weight
sensitivity analyses, and (v) an open-source Python reference
implementation released with a citable DOI is, to the best of our search,
absent from the published literature. The contributions of this paper map
one-to-one onto the five elements of that combination.

% =====================================================================
\section{Model Formulation}\label{sec:model}

\subsection{Problem definition}\label{sec:model-def}

A disaster has just occurred. Civil-protection planners face a bounded
decision: among the $m$ affected people whose locations and conditions are
known, choose at most $n_{\mathrm{dir}}$ to direct to one of $n$ relief centres.
The constraint $n_{\mathrm{dir}} < m$ is operational, not
philosophical: relief capacity in the first hours after a major event is
always smaller than need \cite{cred2023, krishna2021}. The decision-maker does
not get to serve everyone; they have to decide whom to serve, where to send
them, and accept that some will wait.

Three concerns shape every direction decision. The first is who: a moderately
injured 75-year-old living alone has different priority than a healthy adult
with a support network. The second is how: a partially blocked road with
moderate hazard delivers a person but at risk; a longer clear road delivers
the same person more slowly but more safely. The third is where: a centre
already at 80\,\% occupancy will exhaust supplies faster, regardless of how
the route looks. Real practitioners weigh all three simultaneously
\cite{behl2019}.

We model the decision as a multi-objective allocation problem. Each candidate
solution is a partial mapping $d: P \to C \cup \{\bot\}$ that selects exactly
$n_{\mathrm{dir}}$ people from $P$ and assigns each to a centre in $C$. The
remaining $m - n_{\mathrm{dir}}$ people are not directed in this round. The
mapping is feasible by construction; the question is which feasible mappings
sit on the Pareto front under the four objectives we develop below.

\subsection{Resilience-anchored fuzzy indices}\label{sec:model-indices}

We replace the cost-and-distance objective set common in the
allocation literature \cite{tzeng2007, abounacer2014, gama2016, zhao2017}
with four fuzzy-aggregated indices, each anchored in one component of Bruneau et
al.'s \cite{bruneau2003} 4R supply-chain-resilience framework. Three of the
indices form the baseline three-objective formulation; the fourth, and the
splitting of the second, are new here.

A point on terminology, since the distinction matters for what follows. The
\emph{inputs} to these indices are quantitative measurements: a travel duration
in minutes, a distance in kilometres, an occupancy rate, a road-condition
category. What is qualitative is the \emph{output} --- the linguistic resilience
deficit that a fuzzy inference system produces by combining several quantitative
inputs through expert-elicited rules. We therefore call the four objectives
fuzzy-aggregated rather than qualitative: each digests crisp, measurable signals
into a single expert-style judgement on a common $[0,100]$ scale.

\subsubsection{Fairness Level in People Prioritization (FLPP) — Resourcefulness}

FLPP captures whether the right people are being directed first. The 4R
framework's \emph{Resourcefulness} is the system's capacity to identify
problems, mobilise resources, and prioritise correctly under stress
\cite{bruneau2003, hosseini2016}. In a humanitarian last-mile setting that
translates directly into one question: among the $m$ candidates, did we send
the most vulnerable? FLPP encodes a need-weighted Rawlsian reading of
fairness --- prioritisation by greatest individual need rather than equal
shares or proportional allocation. We name the index
``Fairness Level in People Prioritization''; readers
preferring the SCM literature's vocabulary may read it as a Rawlsian
maximin term \cite{soghrati2022}.

We compute FLPP from three sub-indices: a vulnerability score $\mathrm{VS}_j \in
[0,1]$ aggregating age, disability, injury, and living-alone status; an
infrastructure damage level $\mathrm{IDL}_j \in [0,100]$ at the person's
location; and a resource-time-remaining estimate $\mathrm{RTR}_j$ in hours.
Higher VS and IDL with lower RTR mean a person should be directed first; if
they are not, FLPP rises. FIS$_1$ digests the three sub-indices into the
unfairness output $\mathrm{ULPP}_j \in [0,100]$ via a 27-rule rule base
\cite{mamdani1975}.

\subsubsection{Transportation Robustness (TR) — new}

The baseline formulation folds transport robustness and rapidity together into a
single Transportation Infeasibility Level (TIL). For supply-chain resilience
analysis that fusion is too coarse: \emph{Robustness} (the network's capacity
to deliver despite damage) and \emph{Rapidity} (the network's capacity to
deliver in time) are distinct 4R components and the practitioner needs to see
them separately. We split them.

TR captures network damage along the route. The relevant sub-indices are the
road condition score $\mathrm{RCS}_{j,i}$ (clear, partially blocked, or
blocked) and the possible-hazard score $\mathrm{PHS}_{j,i}$ (five-level
discrete categorical from no immediate threat to extreme hazard). FIS$_{2a}$
maps the two scores into a Transport Robustness Deficit
$\mathrm{TRD}_{j,i} \in [0,100]$ with a nine-rule base. Rules cover the full
$3 \times 3$ Cartesian: a clear road with no hazard is very-low TRD; a
blocked road with extreme hazard is very-high TRD.

\subsubsection{Rapidity (RP) — new}

RP captures the speed dimension that TR deliberately excludes. Rapidity in
the 4R framework is context-free: how fast the response delivers, independent
of who is being delivered to or what they are being delivered through
\cite{bruneau2003}. We model it as a single-input fuzzy classifier on travel
duration $\mathrm{TD}_{j,i}$ in minutes, producing a Rapidity Deficit
$\mathrm{RPD}_{j,i} \in [0,100]$ via three rules: short trip $\to$ low RPD,
moderate $\to$ moderate, long $\to$ high.

We considered adding RTR as a second input to FIS$_{2b}$ and rejected it.
Resource time remaining already enters FIS$_1$ to set ULPP; reusing it in
FIS$_{2b}$ would double-count Resourcefulness inside Rapidity and break the
4R-objective bijection that motivates the whole formulation. Rapidity is, and
remains, a property of the route, not of the person.

\subsubsection{Centre Allocation Balance Level (CABL) — Redundancy}

CABL captures whether the load is spread or concentrated. \emph{Redundancy} is
the substitutability of system elements under stress: a network where every
relief centre carries equal load fails gracefully when one centre saturates;
a network where one centre carries everything fails catastrophically
\cite{bruneau2003, christopher2004}. CABL operationalises Redundancy as the
mean imbalance of allocations across centres, weighted by each centre's
current state.

The sub-indices are the centre occupancy rate $\mathrm{COR}_i \in [0,100]$,
the resource depletion rate $\mathrm{RDR}_i \in [0,100]$ at the centre, and
travel duration $\mathrm{TD}_{j,i}$ to the centre. The intuition encoded in
the 27-rule FIS$_3$ rule base is monotone in the obvious direction: a centre
that is already full, depleting fast, and far from the directed person
contributes high CAIL; an empty centre with slow depletion and short travel
contributes low CAIL.

\subsection{Notations and variables}\label{sec:model-notation}

Table~\ref{tab:notation} summarises the notation used throughout the model.
Every notation carried over from the baseline three-objective formulation is
preserved; the additions for the four-objective formulation are flagged in the
rightmost column.

\begin{table}[!t]
\centering
\caption{Notation. Symbols carried over from the baseline three-objective
formulation are unchanged; new symbols for the four-objective formulation are marked.}
\label{tab:notation}
\small
\begin{tabularx}{\textwidth}{l X l}
\toprule
\textbf{Symbol} & \textbf{Description} & \textbf{Origin} \\
\midrule
$P$, $m$ & Set of affected people; $m = |P|$. & baseline \\
$C$, $n$ & Set of relief centres; $n = |C|$. & baseline \\
$n_{\mathrm{dir}}$ & Number of people we direct ($0 < n_{\mathrm{dir}} < m$). & baseline \\
$d_{j,i}$ & Binary indicator: $d_{j,i} = 1$ iff person $j$ directed to centre $i$. & baseline \\
$\mathrm{VS}_j \in [0,1]$ & Vulnerability score of person $j$. & baseline \\
$\mathrm{IDL}_j \in [0,100]$ & Infrastructure damage at person $j$'s location. & baseline \\
$\mathrm{RTR}_j$ (hours) & Resource time remaining for person $j$. & baseline \\
$\mathrm{COR}_i \in [0,100]$ & Centre $i$ occupancy rate. & baseline \\
$\mathrm{RDR}_i \in [0,100]$ & Centre $i$ resource depletion rate. & baseline \\
$\mathrm{TD}_{j,i}$ (min) & Travel duration $j \to i$. & baseline \\
$\mathrm{RCS}_{j,i}$ & Road condition $\in \{0.1, 0.5, 1.0\}$. & baseline \\
$\mathrm{PHS}_{j,i}$ & Possible-hazard score $\in \{0.05, 0.25, 0.5, 0.75, 1.0\}$. & baseline \\
$\mathrm{RWS}_{j,i} \in [0,1]$ & Roadworthiness (high = safe), Eq.~\eqref{eq:rws}. & baseline \\
$\mathrm{ULPP}_j \in [0,100]$ & Unfairness in prioritising person $j$. & baseline \\
$\mathrm{TIL}_{j,i} \in [0,100]$ & Transport infeasibility (3-obj only). & baseline \\
$\mathrm{TRD}_{j,i} \in [0,100]$ & Transport Robustness Deficit. & \textbf{new} \\
$\mathrm{RPD}_{j,i} \in [0,100]$ & Rapidity Deficit. & \textbf{new} \\
$\mathrm{CAIL}_{i,td}$ & Centre allocation imbalance contribution. & baseline \\
$\overline{\mathrm{ULPP}}, \overline{\mathrm{TRD}}, \overline{\mathrm{RPD}}, \overline{\mathrm{CAIL}}$
& Mean of corresponding index over the directed set. & \textbf{4 of 4 new in 4-obj} \\
\bottomrule
\end{tabularx}
\end{table}

The vulnerability score $\mathrm{VS}_j$ is derived from age score
$\mathrm{AS}_j$, disability $\mathrm{DS}_j$, injury $\mathrm{IL}_j$, and
living-status $\mathrm{LS}_j$ via the weighted average
\begin{equation}
    \mathrm{VS}_j = \frac{\mathrm{WAS}\cdot\mathrm{AS}_j + \mathrm{WDS}\cdot\mathrm{DS}_j +
                          \mathrm{WIL}\cdot\mathrm{IL}_j + \mathrm{WLS}\cdot\mathrm{LS}_j}
                         {\mathrm{WAS} + \mathrm{WDS} + \mathrm{WIL} + \mathrm{WLS}},
    \label{eq:vs}
\end{equation}
with the four weights tunable per disaster context (defaults all $1$). The
roadworthiness score is given by
\begin{equation}
    \mathrm{RWS}_{j,i} = 1 - \frac{\mathrm{WRC}\cdot\mathrm{RCS}_{j,i} + \mathrm{WPH}\cdot\mathrm{PHS}_{j,i}}
                                  {\mathrm{WRC} + \mathrm{WPH}}.
    \label{eq:rws}
\end{equation}

\paragraph{Note on Eq.~\eqref{eq:rws}: orientation of RWS.}
We define RWS so that it increases with route quality: the weighted average
of RCS and PHS rises with worse conditions (Blocked $= 1.0$, Extreme $= 1.0$),
so Eq.~\eqref{eq:rws} subtracts it from one. This gives $\mathrm{RWS} = 1$ for
a perfectly safe and clear route and $\mathrm{RWS} = 0$ for an impassable one,
so that $\mathrm{TIL}$ is low exactly when the route is safe and fast.

\subsection{Objective functions}\label{sec:model-objectives}

In words before symbols: each objective averages one resilience deficit over
the people we direct, and the optimiser tries to push all four averages down at
once. The first, $\overline{\mathrm{ULPP}}$, is the mean unfairness of the
prioritisation --- low when the most vulnerable were directed first. The second
and third, $\overline{\mathrm{TRD}}$ and $\overline{\mathrm{RPD}}$, are the mean
robustness deficit of the chosen routes (how exposed they are to damage and
hazard) and the mean rapidity deficit (how slow they are). The fourth,
$\overline{\mathrm{CAIL}}$, is the mean imbalance the decision induces across
centres --- low when load is spread rather than piled onto already-stressed
sites. Each is a number on the same $[0,100]$ scale where lower is better.

These four objectives conflict, which is why a single best solution does not
exist and a Pareto front does. Directing the most vulnerable people (low ULPP)
often means reaching those who live where infrastructure is most damaged, along
the worst routes (high TRD) or the longest ones (high RPD). Choosing the fastest
route (low RPD) can mean accepting a more hazardous one (high TRD), and the
safest route is rarely the quickest. Sending many vulnerable people to the
best-equipped centre lowers ULPP but concentrates load and raises CAIL. No
allocation minimises all four at once; the model's job is to surface the
trade-off surface so the practitioner can choose where on it to sit.

Formally, the model minimises four mean-objective functions over the directed set:
\begin{align}
    \min \overline{\mathrm{ULPP}} &= \frac{\sum_{j=1}^{m} \sum_{i=1}^{n} d_{j,i}\, \mathrm{ULPP}_j}
                                          {n_{\mathrm{dir}}}, \label{eq:obj-ulpp} \\
    \min \overline{\mathrm{TRD}} &= \frac{\sum_{j=1}^{m} \sum_{i=1}^{n} d_{j,i}\, \mathrm{TRD}_{j,i}}
                                          {n_{\mathrm{dir}}}, \label{eq:obj-trd} \\
    \min \overline{\mathrm{RPD}} &= \frac{\sum_{j=1}^{m} \sum_{i=1}^{n} d_{j,i}\, \mathrm{RPD}_{j,i}}
                                          {n_{\mathrm{dir}}}, \label{eq:obj-rpd} \\
    \min \overline{\mathrm{CAIL}} &= \frac{\sum_{i=1}^{n} \mathrm{TCAIL}_i}
                                          {n_{\mathrm{dir}}},   \label{eq:obj-cail}
\end{align}
where $\mathrm{TCAIL}_i = \sum_{j=1}^{m} d_{j,i}\, \mathrm{CAIL}_{i,\mathrm{TD}_{j,i}}$
collects the per-allocation CAIL contributions at centre $i$. The model
constraints are:
\begin{equation}
    \sum_{i=1}^{n} d_{j,i} \le 1\ \forall j, \qquad
    \sum_{j=1}^{m} \sum_{i=1}^{n} d_{j,i} = n_{\mathrm{dir}}, \qquad
    n_{\mathrm{dir}} < m,\ n_{\mathrm{dir}} > 0,\ n > 0,
\end{equation}
i.e., each person is directed to at most one centre, exactly
$n_{\mathrm{dir}}$ are directed, and the directed set is a strict subset of
the affected population.

For the baseline three-objective formulation (used in
Section~\ref{sec:tradeoff-analysis} as the H1 comparator), Eqs.
\eqref{eq:obj-trd} and \eqref{eq:obj-rpd} are replaced by the single
$\overline{\mathrm{TIL}}$ objective using the fused FIS$_2$ pathway of the
baseline formulation. Everything else is identical.

\subsection{Resilience-theoretic mapping (new)}\label{sec:resilience-mapping}

The 4R framework earns its place by doing analytical work, not by labelling.
It fixes what each objective is meant to capture and tells the practitioner how
to read the Pareto front. Table~\ref{tab:4r-mapping} states the bijection.

\begin{table}[!t]
\centering
\caption{Bijection between the four objectives and Bruneau et al.'s 4R
supply-chain-resilience framework \cite{bruneau2003}. Each objective
operationalises exactly one 4R component; the practitioner reading the
Pareto front sees four resilience axes simultaneously.}
\label{tab:4r-mapping}
\small
\begin{tabular}{l l l l}
\toprule
\textbf{4R component} & \textbf{Objective} & \textbf{FIS} & \textbf{Practitioner reading} \\
\midrule
Resourcefulness  & $\overline{\mathrm{ULPP}}$ & FIS$_1$        & ``Did we direct the most vulnerable first?'' \\
Robustness       & $\overline{\mathrm{TRD}}$  & FIS$_{2a}$     & ``Will the routes hold under damage?'' \\
Rapidity         & $\overline{\mathrm{RPD}}$  & FIS$_{2b}$     & ``Will the people get there in time?'' \\
Redundancy       & $\overline{\mathrm{CAIL}}$ & FIS$_3$        & ``Is the load spread or concentrated?'' \\
\bottomrule
\end{tabular}
\end{table}

\subsubsection*{Per-component derivation}

Table~\ref{tab:4r-mapping} states the bijection in summary form; the
mapping is not a label slapped on existing objectives but a derivation
from each 4R component's definition. We walk through the four pairings.

\paragraph{Resourcefulness $\to$ ULPP via FIS$_1$.} Bruneau et al.\
\cite{bruneau2003} define Resourcefulness as ``the capacity to identify
problems, establish priorities, and mobilise resources when conditions
exist that threaten to disrupt some element [...]\ of the system.'' Three
clauses, each present in our model. \emph{Identifying problems} maps to
the three sub-indicators feeding FIS$_1$ (vulnerability score VS,
infrastructure damage IDL at the person's location, resource time
remaining RTR); these are precisely the diagnostic signals a triage
coordinator would consult. \emph{Establishing priorities} maps to the
27-rule rule base of FIS$_1$, which is itself an explicit articulation
of expert priority judgement. \emph{Mobilising resources} maps to the
directing decision $d_{j,i}$ that the optimisation produces. ULPP is
zero exactly when the coordinator's three diagnostics, the rule base's
priority logic, and the directing decision are mutually consistent ---
i.e., when the system is being Resourceful in Bruneau's sense.

\paragraph{Robustness $\to$ TRD via FIS$_{2a}$.} Bruneau et al.\ define
Robustness as ``the inherent strength or resistance in a system to
withstand external demands without degradation or loss of
functionality.'' For a humanitarian last-mile route the relevant
external demand is infrastructure damage and route hazard, and the
relevant function is delivering the directed person to the centre. Two
sub-indicators capture this: RCS$_{j,i}$ encodes the inherent strength
of the road (clear / partially blocked / blocked), and PHS$_{j,i}$
encodes external threat along the route. FIS$_{2a}$ digests both into
TRD --- the deficit relative to a perfectly robust route. TRD is zero
exactly when the route can deliver the directed person regardless of
external conditions, which is the limit case of Bruneau's robustness.

\paragraph{Rapidity $\to$ RPD via FIS$_{2b}$.} Bruneau et al.\ define
Rapidity as ``the capacity to meet priorities and achieve goals in a
timely manner in order to contain losses and avoid future
disruption.'' This is the only 4R component with explicit time content.
For our model the canonical time measure is travel duration TD$_{j,i}$.
We use TD as the single input to FIS$_{2b}$ deliberately: any
multi-input formulation would introduce non-time content into Rapidity
and break the framework's component-orthogonality. FIS$_{2b}$ emits
RPD, the deficit relative to the limiting case of instantaneous
delivery; RPD is zero exactly when delivery is fast enough to contain
losses, which is Bruneau's Rapidity in the limit.

\paragraph{Redundancy $\to$ CAIL via FIS$_3$.} Bruneau et al.\ define
Redundancy as ``the extent to which elements [...]\ are substitutable,
that is, capable of satisfying functional requirements in the event
of disruption.'' For a network of relief centres, substitutability is
the property that any one centre saturating does not prevent the
others from absorbing displaced load. Three sub-indicators capture
the dynamics: COR$_i$ (current occupancy of centre $i$), RDR$_i$
(rate at which centre $i$'s resources deplete), and TD$_{j,i}$
(arrival time, which influences when the directed load reaches the
centre relative to depletion). FIS$_3$ digests these into CAIL, the
imbalance the directing decision induces. Low CAIL means the
directing decision spreads load across centres in a way that
preserves substitutability for the next decision round; high CAIL
means the decision concentrates load on already-stressed centres,
reducing future Redundancy.

\paragraph{Why the derivation matters.}
The bijection acts as a design constraint rather than a post-hoc gloss: each
FIS's input set, rule base, and output range must be explicable in the
vocabulary of exactly one 4R component, and each 4R component must be
operationalised by exactly one FIS. The derivation above shows that constraint
holding across all four pairings. Three theoretical consequences follow, which
later sections draw on.

\bigskip
The bijection has three theoretical consequences worth stating explicitly.

\paragraph{Translation, not aggregation.} A practitioner who plans in 4R
vocabulary can read the Pareto front directly, without translation through a
cost/distance/equity intermediary. Asking ``what is the cheapest solution
that keeps Robustness above 70?'' becomes well-formed and answerable. Asking
``which solution maximises Rapidity at the expense of Resourcefulness?''
becomes well-formed and answerable. The 3-objective formulation, where
Robustness and Rapidity are fused into TIL, makes the second question
unanswerable by construction. We test this empirically in
Section~\ref{sec:tradeoff-analysis}.

\paragraph{Theoretical hygiene.} If two of our objectives mapped to the same
4R component the framework would tell us so by complaining. The bijection
constrains the modeller. Adding RTR to FIS$_{2b}$ would have produced a
better-fitting Rapidity index empirically but would have leaked
Resourcefulness into Rapidity, which the 4R framework forbids, so we leave
RTR out.

\paragraph{What the framework does not do.} The 4R framework does not
prescribe weights between the four components. Bruneau et al.
\cite{bruneau2003} are explicit on this point: a community planner balancing
robustness against rapidity makes a different choice than a triage
coordinator on the second day of a flood. The 4R framework provides the
\emph{vocabulary} for the trade-off, not the trade-off itself. Our model
preserves that property by emitting a four-dimensional Pareto front rather
than a single weighted score; the decision-maker selects the trade-off they
want, in 4R terms, post-hoc.

% =====================================================================
\section{Algorithm Design}\label{sec:algorithms}

\subsection{Fuzzy inference in brief}\label{sec:algo-fis-primer}

Each of our four objectives is produced by a fuzzy inference system (FIS), so
we first describe what a FIS is and why we use one, assuming no prior
familiarity. A FIS is a way of turning numbers into a graded expert judgement
using if--then rules written in words. It suits this problem because the people
who run relief operations reason in linguistic categories. A road is
``partially blocked,'' supplies are ``depleting fast,'' a person is ``highly
vulnerable.'' Forcing those judgements into a single tuned formula throws
away information the practitioner actually has.

A FIS works in three steps. \emph{Fuzzification} reads each crisp input and asks
how strongly it belongs to each linguistic level. The levels are defined by
\emph{membership functions}: simple shapes over the input range that return a
degree of belonging between $0$ and $1$. We use two standard shapes. A
\emph{triangular} membership function, written $\mathrm{tri}[a,b,c]$, rises
linearly from $0$ at $a$ to $1$ at the peak $b$ and falls back to $0$ at $c$; it
models a level with a single most-typical value. A \emph{trapezoidal} function,
$\mathrm{trap}[a,b,c,d]$, rises from $a$ to $b$, stays at $1$ across the plateau
$[b,c]$, and falls to $0$ at $d$; it models a level with a range of equally
typical values, which is convenient for the ends of a scale. A given input
typically belongs partly to two neighbouring levels at once: a vulnerability
reading near the boundary between ``Medium'' and ``High'' might register, say,
$0.27$ Medium and $0.73$ High rather than snapping to one label.

\emph{Rule firing} then evaluates a small rule base of the form ``if vulnerability
is High and the centre's service is Low, then unfairness is High.'' Every rule
whose antecedent has non-zero membership contributes, weighted by how strongly it
fired, so several rules speak at once. \emph{Defuzzification} collapses those
weighted contributions back into one number: we use the centroid (centre of mass)
of the aggregated output, which gives a smooth, monotone response on the common
$[0,100]$ output scale. Concretely, a moderately vulnerable person placed at a
nearby but only medium-service centre fuzzifies across the Medium/High levels,
fires a handful of rules, and defuzzifies to a modest unfairness score rather
than a hard ``fair/unfair'' verdict. A fully worked numerical example of all
three steps is given in the companion documentation accompanying the released
code. The remainder of this section specifies the four rule bases and the
membership partitions they use.

\subsection{Fuzzy inference systems}\label{sec:algo-fis}

We use Mamdani inference \cite{mamdani1975} with centroid defuzzification
throughout, for the same reason as earlier work in this line
\cite{rabiei2021, rabiei2023}: linguistic
if--then rules elicited from disaster-response practitioners are easier to
audit and revise than a tuned crisp model, and the Mamdani conjunction over
triangular and trapezoidal membership functions gives the rule base a clear
geometric interpretation. The four-objective formulation uses four control
systems (FIS$_1$ for ULPP, FIS$_{2a}$ for TRD, FIS$_{2b}$ for RPD, FIS$_3$
for CAIL); the baseline three-objective formulation reuses its FIS$_2$ for
TIL and is invoked only as the H1 comparator in
Section~\ref{sec:tradeoff-analysis}.

Each input variable is partitioned into three or five linguistic levels with
trapezoidal endpoints and triangular interiors; the output universe is a
five-level scale (Very Low, Low, Moderate, High, Very High) on $[0, 100]$
shared across all five FIS for cross-objective comparability. Rule-base sizes
are 27 for FIS$_1$ (the full $3^3$ Cartesian product over VS, IDL, RTR),
9 for FIS$_{2a}$ (RCS $\times$ PHS), 3 for FIS$_{2b}$ (TD alone), and 27 for
FIS$_3$ (the full $3^3$ over COR, RDR, TD). Appendix~\ref{app:rules} lists
every rule.

A practical detail matters for performance. We pre-compute every FIS
evaluation once per problem instance: for the four-objective formulation,
$m$ per-person ULPP values and $4 \cdot m \cdot n$ per-pair (TRD, RPD,
CAIL) values, all stored as numpy arrays. Chromosome evaluation in the GA
loop then reduces to fancy-indexing into those arrays --- no scikit-fuzzy
\cite{mamdani1975} calls inside evolution. The cost of pre-computation
on the medium instance (8 centres, 225 people) is roughly 10\,s; the
amortised cost per generation is essentially zero. This is what makes the
540-run experiment matrix in Section~\ref{sec:performance} tractable.

\subsection{Multi-objective evolutionary algorithms}\label{sec:algo-moea}

Three MOEAs share the same population-management loop and differ only in
their survivor-selection operator. We summarise each.

\paragraph{NSGA-II \cite{deb2002}.}
Non-dominated sorting partitions the combined parent--offspring population
into Pareto fronts. Survivors are taken front by front; when a front
overflows the population quota, the remaining seats are filled by
crowding-distance ranking within that front. The crowding-distance
machinery is the workhorse of NSGA-II for two- and three-objective problems
and the reason it dominates the literature in that regime. Recent runtime
analysis \cite{zheng2024nsga2} shows the mechanism degrades as objective
count grows, motivating reference-point methods at higher objective counts.

\paragraph{NRGA \cite{omar2008nrga}.}
NRGA replaces NSGA-II's crowding-distance tie-breaker with uniform random
sampling within the boundary front. The intuition is that crowding distance
encodes a particular notion of diversity (gap between consecutive
solutions in objective space) which may not match the practitioner's notion;
random sampling avoids that prior. The cost is that some Pareto-optimal
solutions are dropped at random, lowering NNS in our experiments by a few
percent. Whether NRGA helps or hurts depends on the problem geometry and the
chromosome encoding --- a sensitivity we revisit in
Section~\ref{sec:perf-algorithms}.

\paragraph{NSGA-III \cite{deb2014nsga3}.}
NSGA-III replaces crowding distance with reference-point-based niching
designed for many-objective problems. A set of structured reference points
$\{r_k\}$ is laid down in the normalised objective simplex; survivors are
chosen to maintain diversity \emph{relative to those reference points}
rather than relative to each other. The mechanism keeps boundary regions
populated and is the modern default for objective counts $\ge 4$. We use
NSGA-III not because we expect it to dominate NSGA-II at four objectives ---
the empirical evidence in the literature is mixed at that boundary --- but
because its CPU cost is independent of population density in front overflow
cases, which we test in Section~\ref{sec:perf-algorithms}.

\subsection{Chromosome encoding}\label{sec:algo-chromosome}

Each individual is a flat list of length $2 n_{\mathrm{dir}}$, read as two
contiguous slices. The first $n_{\mathrm{dir}}$ entries are the
\emph{person slice}: integer person indices, sampled without replacement from
$\{0, \ldots, m-1\}$, i.e.\ a partial permutation that names exactly which
$n_{\mathrm{dir}}$ of the $m$ candidates are directed in this round. The second
$n_{\mathrm{dir}}$ entries are the \emph{centre slice}: gene $k$ holds the centre
to which the $k$-th directed person is sent. We store each centre gene as a real
value $r_k \in [0, 1)$ and decode it to a centre index by
$\lfloor r_k \cdot n \rfloor$ (clamped to $n-1$ at the upper endpoint). The real
value is therefore just a mutation-friendly stand-in for an integer in
$\{0, \ldots, n-1\}$; the centre slice carries exactly a per-person integer
centre assignment, no more. Every chromosome the operators can produce is
feasible by construction: the person slice is a permutation (no duplicate or
out-of-range person), and every centre gene decodes to a valid centre, so no
repair against infeasibility is ever needed.

A reviewer of the earlier model in this line asked, reasonably, why not encode an
allocation as a single length-$n$ vector with one integer centre per person. That
encoding is the natural choice when \emph{every} person is served. Our problem
directs a strict subset ($n_{\mathrm{dir}} < m$; relief capacity is smaller than
need), so the chromosome must encode two decisions, not one: \emph{which} people
are directed and \emph{where} each goes. The person slice carries the first
decision and the centre slice the second. Restricted to the directed people, the
centre slice \emph{is} exactly the one-integer-per-person encoding that reviewer
suggested; the person slice is the additional structure the subset-selection
problem requires.

The order of the (person, centre) pairs is not semantically meaningful: the
decoder treats them as a set of pairs, not a sequence. Variation operators
preserve only the partial-permutation invariant on the person slice. We use
uniform crossover with repair (children inherit from each parent at each
position with probability $0.5$, then duplicate person indices are replaced
with random unused indices) and per-gene mutation (replacement with an
unused index for the person slice; Gaussian perturbation $\mathcal{N}(0, 0.1)$
clipped to $[0, 1)$ for the centre slice). DEAP's built-in operators do not
fit a partial-permutation + real-valued mixed encoding directly, so we
implement the two operators ourselves; the implementation is open-source.

\subsection{NSGA-III reference points (new)}\label{sec:algo-refpoints}

The reference-point set determines NSGA-III's selection geometry. We use
Das--Dennis structured points with $p = 4$ divisions, yielding
$\binom{p + M - 1}{M - 1} = \binom{7}{3} = 35$ points uniformly distributed
over the four-dimensional objective simplex \cite{deb2014nsga3}. The choice
is conventional and easy to reproduce; we considered the two-layer scheme
that adds points along the simplex boundary for better edge coverage and
deferred it to a single supplementary ablation, to be run only if boundary
diversity proved problematic in the algorithm comparison (it did not; see
Section~\ref{sec:perf-algorithms}).

The 35 reference points are recomputed once at solver start and held fixed
across generations. With population size $100 \ge 35$ the reference set
never overflows and the niching procedure runs without partial-front
corner cases. In the four-objective formulation $p = 4$ produces reference
spacing of $1/4$ along each axis, which we found empirically sufficient for
the Pareto-front densities we observe (typically 100 distinct solutions
per front, comfortably above the reference-point count).

% =====================================================================
\section{Performance Analysis}\label{sec:performance}

\subsection{Experimental setup}\label{sec:perf-setup}

We test three problem sizes drawn at a single seed from a deterministic
generator: small (5 centres, 150 affected people, $n_{\mathrm{dir}} = 50$),
medium (8/225/75), and large (10/300/100). Small and large match the baseline
test instances; medium is new and added to give scaling
analysis a third anchor. Each problem instance is materialised as three CSV
files (people, centres, travel) and is reproducible bit-exact from the seed
via the open-source instance generator at
\texttt{experiments/generate\_instances.py}.

For every cell in the (size $\times$ formulation $\times$ algorithm) matrix
we run 30 repetitions with rep-specific seeds drawn at
$\mathrm{seed}_r = 1000 + 7919 r$. Population size is 100, generations 200,
crossover probability $0.7$, mutation probability $0.2$ --- baseline
hyper-parameters held fixed across the matrix so the comparison is on
algorithm and formulation, not hyper-parameters. NSGA-III uses the
Das--Dennis $p = 4$ reference set described in
Section~\ref{sec:algo-refpoints}. Implementation is the open-source Python
\texttt{presidio\_vol\_assign.allocation} module
(scikit-fuzzy + DEAP + pymoo); experiments ran on a single MacBook Pro
laptop. The full 540-run matrix completed in 13.1 minutes wall time
including six FIS pre-computes (one per (size, formulation) cell) and
540 GA evolutions.

For each metric we test normality with Shapiro--Wilk; the parametric
independent-samples $t$-test is used when both groups admit it and the
non-parametric Mann--Whitney $U$ otherwise (the Mann--Whitney $U$ test is
the Wilcoxon rank-sum test, the standard significance test in the MOEA
benchmarking literature). We report $p$-values under
two-sided alternatives; threshold $\alpha = 0.05$.

We report hypervolume (HV) as the primary Pareto-front quality indicator
throughout, because it captures both convergence and diversity in a single
reference-anchored measure and is the modern default for many-objective
comparison. We also record number of non-dominated solutions (NNS), spacing
(SM), and CPU time. Mean Ideal Distance (MID) is reported only as a diagnostic,
shown last in every table; we do not rank algorithms by it, for the reason
given in the note to Table~\ref{tab:stats-summary}.

\paragraph{Selector validation on DTLZ2.}
Before running the experiment matrix we validate that DEAP's
\texttt{selNSGA3} \cite{deb2014nsga3} functions correctly inside our
solver framework, independent of our humanitarian-allocation encoding.
We run 30 repetitions of NSGA-II and NSGA-III on the standard DTLZ2
benchmark with $M=4$ objectives, $n_{\mathrm{var}}=10$ decision
variables, and DEAP's published continuous variation operators
(\texttt{cxSimulatedBinaryBounded}, \texttt{mutPolynomialBounded}). NSGA-III
achieves $\mathrm{HV} = 0.970 \pm 0.006$ against the standard
$(1.1, 1.1, 1.1, 1.1)$ reference point; NSGA-II achieves $0.797 \pm 0.032$
($p < 10^{-15}$, Mann--Whitney $U$). The literature-expected ordering
(NSGA-III ahead of NSGA-II at $M=4$) holds in our framework; the H2a
result reported below is therefore a property of the humanitarian
allocation problem, not an artifact of selector or framework
configuration.

\subsection{Pareto-front overview}\label{sec:perf-pareto}

Figure~\ref{fig:pareto-fronts} shows one representative four-objective
Pareto front per (size, algorithm) cell --- the rep-0 front, with each line
a Pareto-optimal solution traced across the four objective axes. Two
features are immediately legible. First, the fronts span the full output
universe along $\overline{\mathrm{ULPP}}$ and $\overline{\mathrm{CAIL}}$ but
concentrate in the middle of the $\overline{\mathrm{TRD}}$ and
$\overline{\mathrm{RPD}}$ scales --- a consequence of the synthetic
travel-info distribution we use, where extreme road conditions and very
long travel durations are rare. Second, the three algorithms produce
visually similar front shapes; the differences are quantitative, not
qualitative.

\begin{figure}[!t]
    \centering
    \includegraphics[width=\textwidth]{fig5_pareto_fronts.pdf}
    \caption{Representative 4-objective Pareto fronts (rep 0) for each problem
    size, displayed as parallel coordinates over the four mean objectives:
    $\overline{\mathrm{ULPP}}$ (Mean Unfairness in People Prioritization,
    operationalising Resourcefulness), $\overline{\mathrm{TRD}}$ (Mean Transport
    Robustness Deficit, Robustness), $\overline{\mathrm{RPD}}$ (Mean Rapidity
    Deficit, Rapidity), and $\overline{\mathrm{CAIL}}$ (Mean Center Allocation
    Imbalance Level, Redundancy). Each line is one Pareto-optimal solution;
    every axis is on the FIS output universe $[0,100]$ where lower is better.
    Algorithms are colour-coded; line opacity is reduced to make front
    overlap visible.}
    \label{fig:pareto-fronts}
\end{figure}

\subsection{Algorithm comparison (H2, H4)}\label{sec:perf-algorithms}

% Figure~\ref{fig:algorithm-compare} reports H2 with statistical tests.
\begin{figure}[!t]
    \centering
    \includegraphics[width=\textwidth]{fig7_h2_algorithm_compare.pdf}
    \caption{Hypervolume (a) and CPU time (b) per problem size, mean
    $\pm$ s.d.\ over 30 reps. NSGA-II achieves significantly higher HV
    than NSGA-III at every size ($p<0.001$, Mann--Whitney U); NSGA-III is
    significantly faster than NSGA-II at every size ($p<0.0001$).}
    \label{fig:algorithm-compare}
\end{figure}

Three findings emerge from the panels in
Figure~\ref{fig:algorithm-compare}.

\paragraph{H2a refuted: NSGA-II beats NSGA-III on hypervolume.}
At every problem size, NSGA-II achieves significantly higher mean HV than
NSGA-III on the four-objective formulation. On the large instance the means
are NSGA-II $13.98 \times 10^6 \pm 0.29$ vs. NSGA-III $13.59 \times 10^6 \pm 0.33$
($p < 0.0001$, Mann--Whitney $U$); on medium $18.81 \pm 0.40$ vs.
$18.33 \pm 0.69$ ($p = 0.001$); on small $21.08 \pm 0.47$ vs.
$20.74 \pm 0.62$ ($p = 0.027$). NSGA-III still beats NRGA at every size
($p < 10^{-6}$). The original H2a expectation --- that NSGA-III would
dominate at the largest size --- does not hold for our four-objective
formulation in the encoding we use.

The reason is consistent with the recent runtime analysis of Zheng and Doerr
\cite{zheng2024nsga2}: NSGA-II's degradation as objective count grows is
well-documented but the inflection is not at four objectives. Reference-point
methods like NSGA-III pay off most clearly at five or more objectives, where
crowding distance becomes information-poor. At four objectives we sit at
the boundary --- close enough that NSGA-II's crowding-distance machinery
remains competitive, far enough that NSGA-III is not punished. The
practitioner reading is that NSGA-III's advantage on this problem class is
not Pareto-front quality but something else, which the next finding makes
explicit.

\paragraph{H2c confirmed strongly: NSGA-III is the fastest algorithm.}
NSGA-III is significantly faster than NSGA-II at every problem size:
$0.91 \pm 0.07$\,s vs. $1.45 \pm 0.07$\,s on small ($p < 10^{-6}$),
$1.15 \pm 0.06$\,s vs. $1.69 \pm 0.06$\,s on medium ($p < 10^{-6}$),
$1.45 \pm 0.05$\,s vs. $1.99 \pm 0.05$\,s on large ($p < 10^{-6}$). The
relative speedup is $37\,\%$, $32\,\%$, and $27\,\%$ at small, medium, and
large respectively (mean $32\,\%$); it remains stable as the problem grows. NRGA sits between the two on CPU time
(small 1.16\,s, medium 1.40\,s, large 1.69\,s).

The mechanism is straightforward. NSGA-II's crowding-distance computation
in the boundary front is $O(M N \log N)$ for $M$ objectives and $N$
individuals on that front, dominated by the per-objective sort. NSGA-III's
reference-point niching is $O(M N + M K)$ for $K$ reference points
(here $K = 35$) and avoids the per-objective sort entirely. At four
objectives and 100 individuals the constant-factor difference shows up as
roughly half a second per generation pair, which compounds across 200
generations into the gap we measure. The same gap will persist on larger
problems --- the $O(\log N)$ term in NSGA-II grows; the niching term in
NSGA-III grows linearly. For practitioners running real-time response-phase
recommendations, NSGA-III's speedup translates directly to time-to-decision.

\paragraph{Reference-point sensitivity of H2a.}
The HV in Table~\ref{tab:stats-summary} is computed against the fixed
$(100, 100, 100, 100)$ reference point at the FIS output-universe maxima.
A demanding reviewer will note that an instance-aware reference ---
the per-objective worst-case across the union of all algorithms' fronts
on the same instance, with a 5\,\% margin --- is more discriminating
between near-optimal fronts and is the more conservative methodological
choice. We recompute HV under this scheme. The ranking shifts:
NSGA-III leads NSGA-II at every size on instance-aware HV (small
$59{,}360$ vs $50{,}921$, $p = 0.054$; medium $20{,}960$ vs $19{,}157$,
$p = 0.57$; large $13{,}003$ vs $11{,}916$, $p = 0.20$). None of these
differences is statistically significant at $\alpha = 0.05$, but the
direction is opposite to the fixed-reference ranking. The honest
reading is that the H2a verdict depends on HV reference choice: under
loose reference NSGA-II is significantly ahead, under tight reference
NSGA-III is ahead but not significantly so. The two algorithms are
HV-equivalent in any practitioner-relevant sense; what differs is the
fixed-reference noise floor. The CPU-time finding (H2c) below is
reference-independent and is the robust algorithmic claim of this
paper.

\paragraph{Population-size ablation.}
A natural concern about the H2a refutation is that population size 100
may be too small for NSGA-III's reference-point niching to express its
design advantage --- 35 reference points facing a 100-individual
front leaves the niching machinery underexercised. We tested this
directly: 30 additional repetitions per algorithm on the medium 4-obj
instance with pop\_size doubled to 200. NSGA-II remains significantly
ahead of NSGA-III on HV ($19.66 \times 10^6 \pm 0.39$ vs.
$19.10 \times 10^6 \pm 0.58$, $p < 0.001$, Mann--Whitney $U$);
NSGA-III's CPU advantage actually widens (NSGA-III $2.37 \pm 0.06$\,s
vs. NSGA-II $4.82 \pm 0.05$\,s, a 51\,\% speedup vs.\ the 32\,\%
average at pop = 100). The ranking is robust to population size and
the boundary-objective-count interpretation in the previous paragraph
holds; doubling the population does not flip the H2a verdict.

\paragraph{H4 refuted (inverted): NRGA is faster than NSGA-II on large in
our encoding.} An earlier MATLAB implementation of this model in our group
reported NSGA-II faster than NRGA on the large instance. Our Python
implementation finds the opposite.
On large, NRGA achieves $1.69 \pm 0.05$\,s versus NSGA-II $1.99 \pm 0.05$\,s
($p < 10^{-4}$, Mann--Whitney $U$); on the three-objective formulation the
inversion is the same (NRGA $1.50$\,s vs. NSGA-II $1.79$\,s,
$p < 10^{-4}$). We do not claim to know the cause. The earlier
implementation uses MATLAB R2024b with a two-row matrix encoding;
we implement the model in Python (\texttt{scikit-fuzzy} + \texttt{DEAP}
+ \texttt{pymoo}) with the flattened partial-permutation + real-valued
encoding of Section~\ref{sec:algo-chromosome} and our own variation
operators. Three differences between the implementations could each
explain the inversion: encoding layout, variation-operator behaviour,
and platform/runtime overheads. Without re-implementing the exact
MATLAB pipeline we cannot attribute the inversion to any single one;
both implementations produce Pareto fronts of comparable HV, so the
inversion is not pathological in either direction.

We treat the inversion as a methodological observation, not a contradiction
of the earlier results. The substantive lesson is that NSGA-II-vs-NRGA CPU
rankings on humanitarian allocation problems do not transfer between
implementations without verification. In this paper, the NSGA-II/NRGA
ranking is replaced by the more robust
observation that \emph{NSGA-III is the fastest algorithm on this problem
class in our implementation}, which is a claim we can substantiate from
our own data without depending on a MATLAB-vs-Python comparison.

\subsection{Statistical comparison summary}\label{sec:perf-stats}

Table~\ref{tab:stats-summary} aggregates the means, standard deviations, and
test statistics. Reading the four-objective formulation by the primary metric
first: NSGA-II holds a significant HV lead over NSGA-III at every size, while
NSGA-III is significantly faster on CPU time --- the two binding metrics, moving
in opposite directions. The remaining metrics do not separate the algorithms in
any practically meaningful way: NSGA-II and NSGA-III achieve identical NNS
(100 of 100, i.e.\ every individual on the boundary front), and their SM and MID
differences, though statistically detectable (Mann--Whitney $p < 0.001$,
Appendix~\ref{app:stats}), are negligible in magnitude (means equal to one
decimal place). NRGA loses 0--1\,\% of NNS to its uniform-random
tie-breaker, has the lowest HV at every size, and sits between the other
two on CPU time. The CPU ranking is stable across all three problem sizes and
both formulations; the HV lead of NSGA-II over NSGA-III is significant
throughout the four-objective formulation but, in the three-objective baseline,
only at the large size (Appendix~\ref{app:stats}).

\begin{table}[!t]
\centering
\caption{Mean $\pm$ s.d.\ over 30 reps per (size, algorithm) cell on the
four-objective formulation. HV (in millions) is the primary quality indicator;
CPU time is in seconds. MID is shown last as a diagnostic only (see note). The
$p$-values underlying these means are the Mann--Whitney $U$ (Wilcoxon rank-sum)
tests against NSGA-II in Appendix~\ref{app:stats}.}
\label{tab:stats-summary}
\small
\begin{tabular}{l l r r r r r}
\toprule
Size & Algorithm & NNS & HV ($\times 10^6$) & SM & CPU (s) & MID$^{\dagger}$ \\
\midrule
\multirow{3}{*}{small}
  & NSGA-II  & $100.0$ & $21.08 \pm 0.47$ & $0.81 \pm 0.06$ & $1.45 \pm 0.07$ & $77.7 \pm 0.5$ \\
  & NRGA     & $99.1$  & $20.21 \pm 0.87$ & $0.79 \pm 0.05$ & $1.16 \pm 0.06$ & $77.4 \pm 0.5$ \\
  & NSGA-III & $100.0$ & $20.74 \pm 0.62$ & $0.81 \pm 0.06$ & $0.91 \pm 0.07$ & $77.7 \pm 0.5$ \\
\midrule
\multirow{3}{*}{medium}
  & NSGA-II  & $100.0$ & $18.81 \pm 0.40$ & $0.78 \pm 0.05$ & $1.69 \pm 0.06$ & $69.5 \pm 0.4$ \\
  & NRGA     & $99.9$  & $17.35 \pm 0.55$ & $0.77 \pm 0.04$ & $1.40 \pm 0.05$ & $69.3 \pm 0.4$ \\
  & NSGA-III & $100.0$ & $18.33 \pm 0.69$ & $0.79 \pm 0.05$ & $1.15 \pm 0.06$ & $69.5 \pm 0.4$ \\
\midrule
\multirow{3}{*}{large}
  & NSGA-II  & $100.0$ & $13.98 \pm 0.29$ & $0.81 \pm 0.05$ & $1.99 \pm 0.05$ & $58.6 \pm 0.4$ \\
  & NRGA     & $99.8$  & $12.86 \pm 0.29$ & $0.81 \pm 0.04$ & $1.69 \pm 0.05$ & $58.6 \pm 0.4$ \\
  & NSGA-III & $100.0$ & $13.59 \pm 0.33$ & $0.81 \pm 0.05$ & $1.45 \pm 0.05$ & $58.6 \pm 0.4$ \\
\bottomrule
\end{tabular}

\smallskip
{\footnotesize $^{\dagger}$\,MID (Mean Ideal Distance) is reported as a
diagnostic only and is not used to rank algorithms. MID rewards solutions that
sit close to the ideal point, which penalises legitimate Pareto-optimal
trade-offs: every solution on a front is equally non-dominated, so a
front-quality indicator should not prefer the ones nearest a single corner. We
retain MID for backward comparability with the earlier formulation in this line
of work, but read front quality from HV, which captures convergence and
diversity together.}
\end{table}

The MID and SM values are nearly identical across algorithms within each
size, and the small differences are not statistically significant. The binding
metrics for the H2 verdict are HV and CPU time, where the algorithms differ in
opposite directions; MID and SM add no separating information.

% =====================================================================
\section{Pareto Trade-off Analysis (4-obj vs 3-obj)}\label{sec:tradeoff-analysis}
% TIER 1 PLACEHOLDER — target ~1{,}200 words. Tests H1.

The four-objective formulation has to do empirical work, not just
theoretical work. The H1 hypothesis tests whether splitting the baseline TIL
into separate Robustness (TRD) and Rapidity (RPD) components actually
exposes Pareto trade-offs the three-objective formulation cannot represent.
We test it two ways: rank correlation between the split axes (do they carry
independent information?) and projection dominance (what happens when we
fold the 4-obj front into 3-obj space?).

\subsection{Spearman rank correlation between TRD and RPD}\label{sec:h1-spearman}

If TRD and RPD were redundant --- two names for the same thing along the
Pareto front --- the rank correlation between them would be near unity in
absolute value, and splitting TIL would add nothing. We compute the Spearman
$\rho(\overline{\mathrm{TRD}}, \overline{\mathrm{RPD}})$ across each
four-objective Pareto front in the 270-run cell.
Figure~\ref{fig:h1-spearman}(a) shows the distribution.

The mean correlation across all 270 fronts is $\rho = -0.29$, with $82\,\%$
of fronts under the H1 threshold $|\rho| < 0.5$. The negative sign is itself
informative: along the Pareto front, solutions that achieve low TRD (clear,
safe routes) tend to come at the cost of higher RPD (longer travel times),
and vice versa. The two axes are not redundant; they expose a genuine
trade-off where the practitioner picks a fragile-but-fast route or a
robust-but-slow one. NSGA-III shows the cleanest separation
($93$--$97\,\%$ of fronts under threshold across the three sizes), NSGA-II
is close behind ($73$--$87\,\%$), and NRGA's slightly noisier fronts
produce the lowest pass rate ($70$--$80\,\%$). The mechanism is consistent
with Section~\ref{sec:perf-algorithms}: NRGA's uniform-random tie-breaker
sometimes drops boundary solutions that otherwise extend the trade-off
range.

\subsection{Projection-dominance analysis}\label{sec:h1-projection}

The Spearman result establishes that the two axes carry independent
information. The projection analysis answers a sharper question: how many
4-obj Pareto-optimal solutions \emph{disappear} when we collapse to 3-obj
space using the baseline TIL fusion?

For each 4-obj front, we re-evaluate every (person, centre) allocation
through the baseline FIS$_2$ pathway --- compute the RWS by
Eq.~\eqref{eq:rws}, then evaluate FIS$_2$(TD, RWS) to get TIL. Averaging
gives a per-solution $\overline{\mathrm{TIL}}$, and the projected solution
is $(\overline{\mathrm{ULPP}}, \overline{\mathrm{TIL}}, \overline{\mathrm{CAIL}})$.
We then test Pareto dominance among the 100 projected points: a projected
solution is \emph{lost} if some other projected solution dominates it
strictly in all three dimensions.

Figure~\ref{fig:h1-spearman}(b) shows the result. The mean fraction of
4-obj solutions whose 3-obj projection is dominated is $0.51$ on small,
$0.49$ on medium, and $0.46$ on large --- around half the front, in every
cell, across all three algorithms. \emph{Half} the Pareto-optimal solutions
that the four-objective formulation surfaces are invisible in the
three-objective formulation, in the sense that another projected solution
dominates them along all three TIL-fused axes simultaneously.

The economic-intuition reading: the 3-obj formulation encodes a particular
trade-off between Robustness and Rapidity --- the one defined by the RWS
weighting in Eq.~\eqref{eq:rws} --- and lets only solutions consistent with
that fixed trade-off through. The 4-obj formulation makes the trade-off a
free parameter the practitioner picks at decision time, and so admits
solutions across the entire Robustness/Rapidity plane.

\paragraph{Null-model control: is the loss informational or geometric?}
A fair question follows immediately. \emph{Some} fraction of any 4D
Pareto front will be dominated when projected to 3D simply because
projection loses information --- a geometric effect, not specific to the
baseline RWS-weighted fusion. To separate the geometric component from the
informational, we sample 100 random convex fusions
$\mathrm{TIL}_{\mathrm{null}} = \alpha \cdot \overline{\mathrm{TRD}}
+ (1-\alpha) \cdot \overline{\mathrm{RPD}}$ with
$\alpha \sim \mathcal{U}(0,1)$ per fusion, recompute the projection-
dominance fraction for each, and form an empirical null distribution
against which the baseline-fusion fraction is compared.

The result tempers the 50\% headline. Pooled across all 270 fronts, the
baseline-fusion mean dominance fraction is 0.486; the null-distribution
mean is 0.479. The difference is $+0.008$ --- under one percentage point.
The mean z-score of the baseline fraction in the per-front null distribution
is $+0.09$; only $7.4\%$ of fronts have a baseline fraction above the
null 95th percentile, and $5.6\%$ below the 5th. The projection-dominance
loss is, almost entirely, geometric.

What survives this control is the Spearman finding from
Section~\ref{sec:h1-spearman}: TRD and RPD carry independent rank
information. The 4-obj formulation does expose new trade-offs --- the
Spearman result establishes that --- but the specific 3-obj fusion the
baseline uses is not worse than a random fusion at preserving them; it is
not the source of the dominance loss. We rewrite the H1 finding
accordingly in Section~\ref{sec:h1-verdict}.

\subsection{Practitioner trade-off vignettes}\label{sec:h1-vignettes}

Aggregate statistics tell us that half the 4-obj front is invisible in
3-obj space; they do not tell us \emph{which} half. Three vignettes,
constructed from solutions that appear on the 4-obj front but are
dominated in the 3-obj projection, illustrate what the Robustness/Rapidity
split unlocks for the decision-maker.

\paragraph{Vignette 1 --- the fragile-but-fast route.}
A directed person reaches a relief centre via a route with PHS = significant
hazard but TD = 25\,min. The 3-obj formulation hides this option: TIL fuses
the fragility into the same axis as the speed, and the solution is
dominated in 3-obj space by a slower, safer alternative even when both
options serve the same person. The 4-obj formulation surfaces it as a
distinct point on the Pareto front, low on RPD and high on TRD, letting the
practitioner pick it deliberately if speed is the binding constraint
(``the patient must reach the centre within 30 minutes; we accept the
hazard'').

\paragraph{Vignette 2 --- the slow-but-robust route.}
The mirror case: TD = 90\,min on a clear road with no hazard. The 4-obj
front shows it as low TRD, high RPD --- the conservative choice. The 3-obj
formulation again merges the two axes; the practitioner cannot articulate
``I want robustness even at the cost of an extra hour'' without
unfolding TIL into its components, which the 4-obj formulation does for
them.

\paragraph{Vignette 3 --- the capacity-constrained centre under load.}
A centre at COR = 80\,\%, RDR = 70\,\% (rapid depletion), with TD = 30\,min
to a cluster of vulnerable people. The CAIL objective penalises directing
more people there; the ULPP objective rewards it (the people are vulnerable).
The 4-obj front exposes the Pareto trade-off between Resourcefulness and
Redundancy directly --- send the vulnerable here and accept higher CAIL, or
send them to a less-loaded centre and accept higher ULPP. The 3-obj
formulation has the same Resourcefulness-vs-Redundancy axis (it is
unchanged by the split), but it cannot tell the practitioner whether the
robustness of the route to this centre is also a constraint --- the third
vignette exposes the multi-axis nature of the realistic decision in a way
that fusing TIL obscures.

\subsection{H1 verdict}\label{sec:h1-verdict}

The two halves of the H1 test land asymmetrically and deserve separate
verdicts.

\paragraph{Spearman half: confirmed.} TRD and RPD carry independent rank
information across the 4-obj fronts. Pooled across all 270 4-obj runs,
the mean Spearman correlation is $\rho = -0.29$ and $82\,\%$ of fronts
satisfy $|\rho| < 0.5$. The two axes are not redundant; the split exposes
information the fused TIL cannot represent.

\paragraph{Projection-dominance half: geometric, not informational.}
The 50\% dominance fraction we observe under the baseline RWS fusion is
within one percentage point of the null mean from random convex fusions
of TRD and RPD. The dominance loss is essentially geometric --- a
property of projecting from 4D to 3D, not of the baseline's particular
3-objective formulation. We do not claim that the 3-obj formulation is
specifically destructive; we claim that the 4-obj formulation surfaces
trade-off information that any 3-obj fusion would compress.

\paragraph{Joint reading.} The 4-objective formulation is empirically
informative in the sense that the Robustness and Rapidity axes are
independent (Spearman). It is \emph{not} empirically richer than the
3-obj formulation in the sense of preserving more solutions through
projection (geometric). The practitioner gain is therefore in
\emph{visibility}: the 4-obj front lets the decision-maker articulate
Robustness/Rapidity trade-offs at decision time, whereas the 3-obj
fusion bakes a single trade-off in and exposes only one axis. The 4R
reframing of Section~\ref{sec:resilience-mapping} delivers this
visibility; it is operationally useful even if its dominance footprint
is not statistically distinguishable from random fusion.

\begin{figure}[!t]
    \centering
    \includegraphics[width=\textwidth]{fig6_h1_spearman.pdf}
    \caption{H1 evidence. (a) Distribution of Spearman
    $\rho(\overline{\mathrm{TRD}}, \overline{\mathrm{RPD}})$ across 270 4-objective
    runs; 82\% lie within $|\rho|<0.5$. (b) Mean fraction of solutions whose
    3-objective projection is strictly dominated by another projection in the
    same front; the dashed line marks the 20\% threshold for H1
    confirmation.}
    \label{fig:h1-spearman}
\end{figure}

% =====================================================================
\section{Sensitivity Analysis}\label{sec:sensitivity}
% TIER 1 PLACEHOLDER — target ~1{,}200 words. Tests H3.

\subsection{Rule-base perturbation (H3a)}\label{sec:sens-rules}

\begin{figure}[!t]
    \centering
    \includegraphics[width=\textwidth]{fig8_h3a_rule_sensitivity.pdf}
    \caption{Per-FIS distribution of $\Delta$HV under single-rule deletion
    on the medium 4-obj problem (NSGA-II, 30 reps). The dashed band marks
    $\pm 5\%$, the H3a confirmation criterion. FIS$_1$, FIS$_{2a}$, and FIS$_3$
    have median $|\Delta\mathrm{HV}| < 0.2\%$; FIS$_{2b}$, with only three rules,
    is structurally more sensitive (single-rule deletion equals a
    33\% rule-base reduction).}
    \label{fig:rule-sensitivity}
\end{figure}

A reviewer or practitioner who reads the recommendations of a fuzzy expert
system has a fair question waiting: are the recommendations artifacts of
the specific rule base elicited from this set of experts, or would slightly
different experts produce essentially the same Pareto front? We test it by
deleting one rule at a time and measuring the impact.

For each of the 66 single-rule deletions across the four FIS rule bases
(FIS$_1$: 27 rules, FIS$_{2a}$: 9, FIS$_{2b}$: 3, FIS$_3$: 27), we run
NSGA-II $\times$ 30 reps on the medium 4-obj instance and record the mean
hypervolume against a 30-rep baseline run with the full rule base. The
H3a confirmation criterion is median $|\Delta\mathrm{HV}| \le 5\,\%$ per
FIS \emph{and} worst-case $|\Delta\mathrm{HV}| \le 20\,\%$.
Figure~\ref{fig:rule-sensitivity} shows the per-FIS distributions on
the medium instance. We additionally repeat the sweep on the small and
large instances with 5 reps each, to test whether the verdicts are
size-dependent (Table~\ref{tab:h3a-by-size}).

\begin{table}[!t]
\centering
\caption{Per-FIS median absolute $\Delta\mathrm{HV}$ under single-rule
deletion, by problem size. Confirmation band $\le 5\,\%$. FIS$_{2b}$'s
three-rule structure produces near-threshold sensitivity on small and
medium and below-threshold on large; the larger rule bases are robust
at every size.}
\label{tab:h3a-by-size}
\small
\begin{tabular}{l c c c c}
\toprule
FIS & rules & small (5 reps) & medium (30 reps) & large (5 reps) \\
\midrule
FIS$_1$       & 27 & 1.38\,\% & 0.39\,\% & 1.20\,\% \\
FIS$_{2a}$    & 9  & 0.72\,\% & 1.03\,\% & 1.00\,\% \\
\textbf{FIS$_{2b}$}  & \textbf{3}  & \textbf{5.78\,\%} & \textbf{6.09\,\%} & 2.87\,\% \\
FIS$_3$       & 27 & 0.00\,\% & 0.32\,\% & 1.05\,\% \\
\bottomrule
\end{tabular}
\end{table}

The 27-rule and 9-rule systems are robust across all three sizes: median
$|\Delta\mathrm{HV}|$ stays under $1.4\,\%$ for FIS$_1$, FIS$_{2a}$, and
FIS$_3$. On the medium instance only $3.7\,\%$ of FIS$_1$ single-rule
deletions exceed the $\pm 5\,\%$ band, and only $7.4\,\%$ of FIS$_3$
deletions. These three rule bases confirm H3a comfortably at every
problem size.

\paragraph{The FIS$_{2b}$ caveat is structural, not pathological.}
FIS$_{2b}$ has only three rules --- the single-input Rapidity classifier
mapping travel duration $\in \{$short, moderate, long$\}$ to RPD $\in
\{$low, moderate, high$\}$. Single-rule deletion in a three-rule FIS is
a $33\,\%$ rule-base reduction, not a fine perturbation, and the
sensitivity reflects that: median $|\Delta\mathrm{HV}|$ is $5.78\,\%$
on small, $6.09\,\%$ on medium, and $2.87\,\%$ on large. The first two
narrowly breach the $5\,\%$ confirmation threshold; the worst-case
deletion (rule 0, short $\to$ low) produces $|\Delta\mathrm{HV}|$ up to
$25\,\%$. The criterion is failed in the strict reading for FIS$_{2b}$.

Three observations temper this. First, the breach is structural: the
3-rule FIS would fail any single-rule-deletion stress test by
construction. Second, the breach diminishes with problem size --- the
large instance falls comfortably back inside the band, because larger
$n_{\mathrm{dir}}$ averages over more allocations and smooths
single-rule impact. Third, the breach is confined to FIS$_{2b}$; the
remaining $63$ of $66$ rule deletions across the four FISs stay well
inside the band. The practical implication is a design observation for
future versions: FIS$_{2b}$ should be expanded to a five-level Rapidity
classifier (Very Short, Short, Moderate, Long, Very Long $\to$
five-level RPD), matching the linguistic granularity of the other FISs
and reducing the structural sensitivity of single-rule perturbations.
We leave this as future work and report H3a as \emph{confirmed for the
three larger rule bases at every size; narrowly failed for FIS$_{2b}$
on small and medium for structural reasons traceable to its three-rule
design}.

\subsection{Weight perturbation (H3b)}\label{sec:sens-weights}

\begin{figure}[!t]
    \centering
    \includegraphics[width=\textwidth]{fig9_h3b_weight_sensitivity.pdf}
    \caption{Mean objective values across 100 Latin-Hypercube samples of
    weights $\in$ baseline$\cdot$(1$\pm$20\%) on the medium 4-obj NSGA-II
    instance. Per-objective coefficients of variation are all below 1\% — far
    below the 10\% threshold for H3b confirmation. The dashed line marks the
    sample mean.}
    \label{fig:weight-sensitivity}
\end{figure}

The complementary question is parametric. The vulnerability score VS in
Eq.~\eqref{eq:vs} and the roadworthiness score RWS in Eq.~\eqref{eq:rws}
each depend on weights elicited from experts. If the practitioner consults
a different panel of experts and gets weights $20\,\%$ different in either
direction, do the Pareto-front recommendations change?

We answer it with 100 Latin-Hypercube samples on
$(\mathrm{WAS}, \mathrm{WDS}, \mathrm{WIL}, \mathrm{WLS}, \mathrm{WRC},
\mathrm{WPH})$ drawn from baseline$\cdot(1 \pm 0.20)$, each followed by 30
NSGA-II reps on the medium 4-obj instance. The H3b confirmation criterion
is coefficient of variation $\mathrm{CV} = \sigma / \mu \le 10\,\%$ for each
of the four mean objectives across the 100 samples.
Figure~\ref{fig:weight-sensitivity} shows the per-objective sample
distributions.

The result is far stronger than the threshold demands. The four
coefficients of variation on the medium instance are $0.31\,\%$ (ULPP),
$0.36\,\%$ (TRD), $0.12\,\%$ (RPD), and $0.59\,\%$ (CAIL) --- between
one and two orders of magnitude inside the confirmation band. The
hypervolume CV is $0.47\,\%$; the NNS CV is $0.01\,\%$. Per-sample
mean ULPP varies from 50.84 to 51.67 across the entire $\pm 20\,\%$
weight box; mean RPD varies by less than 0.18 units on a 100-unit
scale. We repeat the sweep with 20 LHS samples on the small and large
instances; CVs remain below $1\,\%$ for every objective at every size
(small: ULPP $0.42\,\%$, TRD $0.52\,\%$, RPD $0.12\,\%$, CAIL
$0.96\,\%$; large: $0.36\,\%$, $0.37\,\%$, $0.10\,\%$, $0.17\,\%$).
H3b is confirmed at every problem size.

The mechanism is the FIS's linguistic discretisation. A person whose
baseline $\mathrm{VS} = 0.62$ falls into the High band of FIS$_1$; a
$\pm 20\,\%$ perturbation of the four VS weights moves $\mathrm{VS}$ within
roughly $\pm 0.05$, which is rarely enough to push them into the Medium
band. The dominant rule firing on that person is unchanged across
practically every sample, and so is the ULPP output. The same logic applies
to RWS feeding the baseline FIS$_2$ pathway, although that pathway is not
used in the 4-obj formulation directly. The aggregate effect on the Pareto
front is that $\pm 20\,\%$ weight perturbations propagate as $\le 1\,\%$
front-mean changes --- the GA finds essentially the same trade-offs
regardless of how the practitioner sets the weights inside this band.

\subsection{Practitioner takeaway}\label{sec:sens-takeaway}

The two sensitivity studies answer different versions of the same question
and converge, with one qualification, on the same answer: the
four-objective Pareto-front recommendations are robust to the kinds of
expert-elicitation noise that real practitioners produce.

For \emph{rule-base} elicitation noise --- ``one of these if-then rules
is wrong but I'm not sure which'' --- the three larger rule bases
(FIS$_1$, FIS$_{2a}$, FIS$_3$) have median $|\Delta\mathrm{HV}|$ under
$1.4\,\%$ at every problem size. The three-rule FIS$_{2b}$ is the
exception: its median $|\Delta\mathrm{HV}|$ reaches $5.8$--$6.1\,\%$ on
the small and medium instances, narrowly outside the confirmation band.
This is a structural property of three-rule classifiers, not a flaw in
the methodology, and it points to a concrete fix (a five-level
FIS$_{2b}$, Section~\ref{sec:disc-limitations}). For \emph{weight}
elicitation noise --- ``my colleagues would weight age slightly more
than disability than I do'' --- the impact on mean objectives is under
$1\,\%$ across the realistic $\pm 20\,\%$ band, at every problem size.

The practitioner reading is that consulting multiple panels of experts is
not a substitute for consulting good ones, but it is also not necessary for
robust recommendations. A single elicitation session with one panel,
followed by the four-objective optimisation we describe, will produce
Pareto fronts indistinguishable from those produced by averaging five
independent elicitations within $\pm 20\,\%$. This makes the methodology
deployable in real response-phase settings where elicitation time is a
binding constraint.

% =====================================================================
\section{Discussion}\label{sec:discussion}

The five hypothesis tests in Sections~\ref{sec:performance}--\ref{sec:sensitivity}
land asymmetrically. H2c and H3b are confirmed cleanly. H1 is confirmed
in its Spearman half (the Robustness and Rapidity axes carry independent
rank information) but its projection-dominance half is geometric rather
than informational. H3a is confirmed for the three larger FIS rule
bases at every problem size but narrowly failed for the three-rule
FIS$_{2b}$ on small and medium instances. H2a is reference-sensitive:
NSGA-II leads on hypervolume under a loose reference point, NSGA-III
under a tight one, with no statistically significant separation under
the tight reference. H4 is refuted in the inverted direction, with the
cause undetermined. We read this asymmetry as a feature of an honest
evaluation rather than a weakness: the qualified verdicts are where the
substance lies. Three threads of implication follow.

\subsection{Findings vs.\ supply-chain-resilience theory}\label{sec:disc-scr}

The H1 confirmation is the substantive theoretical result. Splitting TIL
into separate Robustness and Rapidity objectives, motivated by the 4R
framework, produces a Pareto front with around half its solutions
unrepresentable in the fused 3-objective formulation
(Section~\ref{sec:h1-projection}). The split does practical work: it surfaces a
real Robustness/Rapidity trade-off that practitioners articulate and policy
documents reference, one that disappears under single-axis fusion. This is the first published evidence we are aware of
that explicitly operationalising 4R inside a humanitarian-allocation model
\emph{adds empirical content beyond the theoretical framing}. The
implication for SCR theory is concrete: 4R is not just a description
language for resilience; when it is used to organise the objective set of
a computational allocation model, the resulting front is materially
different from one organised around classical OR objectives.

\subsection{Findings vs.\ MOEA literature}\label{sec:disc-moea}

The H2 split (refuted on HV, confirmed on CPU) sharpens the practitioner
recommendation for NSGA-III on this problem class. NSGA-III is the right
algorithmic choice not because it dominates the Pareto front but because
it produces a \emph{Pareto-front-equivalent recommendation} 32\% faster
on average (27--37\% across sizes) than NSGA-II at every problem size we tested
(Section~\ref{sec:perf-algorithms}). For the response-phase setting where
time-to-decision is the binding operational constraint, the speed
advantage is more useful than a marginal HV advantage would have been.
The finding aligns with the runtime analysis in \cite{zheng2024nsga2}: the
inflection point at which reference-point methods take over from
crowding-distance methods is above four objectives for HV but below four
objectives for CPU efficiency.

The H4 inversion is the methodological reminder. An earlier MATLAB
implementation of this model reported NSGA-II faster than NRGA on the large
instance; we observe the opposite. We do not isolate the cause: that earlier
pipeline (MATLAB R2024b, two-row encoding, MATLAB-default
\texttt{gamultiobj} operators) and ours (Python, flattened
encoding, custom DEAP operators) differ along multiple axes
simultaneously, and we did not re-implement the MATLAB operators to
disentangle them. For the FIS-MOEA literature the lesson stands:
comparative claims about NSGA-II vs.\ NRGA on this problem class do not
transfer between implementations without verification. The more robust
empirical claim within our implementation is that NSGA-III is faster
than both NSGA-II and NRGA, which is reference-independent (see
Section~\ref{sec:perf-algorithms}, reference-point sensitivity
paragraph).

\subsection{Practitioner implications}\label{sec:disc-practitioners}

Three implications follow for the civil-protection planner.

First, the 4R vocabulary maps onto the optimisation output. A planner who
reads policy documents in 4R terms can read the Pareto front in 4R terms;
no translation step intervenes. ``Show me the Pareto-optimal solutions
with Robustness above 70'' becomes well-formed
(Section~\ref{sec:resilience-mapping}). For the SI's audience of
supply-chain resilience researchers and civil-protection planners, this
alignment of optimisation tooling with policy vocabulary is the most
immediately actionable contribution.

Second, the recommendations are robust under realistic
expert-elicitation noise (Section~\ref{sec:sensitivity}). A planner using
the model with a single panel of experts will produce Pareto fronts
indistinguishable from one using five averaged panels within $\pm 20\,\%$
weight perturbation. Single-rule deletion in the larger FISs (FIS$_1$,
FIS$_3$) shifts the front by under $0.2\,\%$ in the median. The planner
does not need to over-invest in elicitation breadth.

Third, the artifact is open. Every result in this paper is reproducible
from \url{https://github.com/presidio-v/presidio-hardened-vol-assign} at tag
\texttt{v0.2.0}, archived on Zenodo (DOI:~\url{https://doi.org/10.5281/zenodo.21083547}).
The Python CLI
runs on a laptop; the medium-instance experiments complete in minutes.
This is not a marginal contribution to the FIS-MOEA disaster line: prior
work in this line \cite{rabiei2021, rabiei2023} uses MATLAB
without releasing implementations, and a reviewer or practitioner who
wanted to check or extend the published work had no entry point.

\subsection{Limitations}\label{sec:disc-limitations}

Three limitations bound the present work.

The instances are synthetic. People, centres, and travel attributes are
drawn from generators with parameters chosen to span the FIS rule space
representatively, not from a documented event. We argue in
Section~\ref{sec:disc-future} for a real-event case study as the next
deliverable; the present results establish the methodology on
controllable instances first.

The FIS rule bases reflect one set of expert-elicited judgements ---
those carried over from the baseline formulation for FIS$_1$ and FIS$_3$ and
extended to FIS$_{2a}$ and FIS$_{2b}$ following the same pattern. The
sensitivity analysis in Section~\ref{sec:sensitivity} establishes that
small perturbations of these judgements do not change recommendations
materially, but it does not establish that a panel of experts from a
different humanitarian-logistics tradition would converge on the same
rule base. The model is robust within elicitation noise; the elicitation
itself is a methodological choice we inherit, not derive.

The Rapidity FIS (FIS$_{2b}$) has only three rules, and
Section~\ref{sec:sens-rules} shows that single-rule deletion in
FIS$_{2b}$ produces median $|\Delta\mathrm{HV}|$ of $5.8\,\%$ and
$6.1\,\%$ on the small and medium instances --- narrowly outside the
$5\,\%$ H3a confirmation band --- with a worst-case deletion near
$25\,\%$. This is structural: removing one of three rules is a $33\,\%$
rule-base reduction, and any three-rule FIS would fail a single-rule-
deletion stress test by construction. The breach is confined to
FIS$_{2b}$, diminishes with problem size (the large instance falls
back inside the band at $2.9\,\%$), and does not touch the three
larger rule bases, which confirm H3a at every size. The honest reading
is that the three-rule single-input classifier sits at the lower bound
of useful FIS granularity; expanding it to a five-level Rapidity
classifier is the indicated future-version fix. We report H3a as a
qualified confirmation rather than a clean one, and treat FIS$_{2b}$'s
structural sensitivity as a documented limitation of the current
model.

\subsection{Future work}\label{sec:disc-future}

Four directions follow naturally.

\paragraph{Explicit-baseline benchmark and hard constraints.}
The comparison in this paper is between solvers on our own model, which
characterises the model but does not yet measure it against an external
allocation method. The planned next extension benchmarks the four-objective
framework against an explicit baseline --- a deterministic weighted-sum greedy
heuristic and an exact weighted-sum optimum (Hungarian assignment and
mixed-integer programming) --- so that the framework's quality is read against a
known reference rather than only against alternative solvers. The same extension
replaces the present soft-capacity treatment with optional hard shelter-capacity
and transport-distance constraints. Both the baseline comparators and the
hard-constraint mode are already implemented in the released tool; reporting them
as a full empirical study is deferred to keep the present contribution focused on
the resilience reframing and its robustness.

\paragraph{Real-event case study.} The next version (v0.3.0 in our
roadmap) applies the four-objective allocation model to a documented
event --- the 2023 T\"urkiye--Syria earthquake or the 2024 Valencia
floods are the candidates we are evaluating --- using EM-DAT,
OpenStreetMap, and INFORM Risk Index data. The empirical hypotheses
remain the same but the validation gains anchoring in observed
practitioner decisions.

\paragraph{Multi-modal transport.} The current Rapidity index assumes a
single transport mode (vehicle on road). Extending FIS$_{2b}$ to take a
mode indicator as a second input --- on-foot, vehicle, public transit,
boat --- would let the model represent multi-modal evacuation routes
without losing the 4R-axis discipline.

\paragraph{Algorithmic alternatives.} MOEA/D and RVEA are the obvious
candidates to add to the comparison \cite{deb2014nsga3}. Both are
reference-direction methods like NSGA-III but with different niching
mechanisms. The H2c finding suggests NSGA-III dominates the relevant
algorithmic axis at four objectives; widening the comparison would
strengthen or weaken that claim.

% =====================================================================
\section{Conclusions}\label{sec:conclusion}

We have presented a humanitarian last-mile allocation model whose four
objectives map one-to-one onto Bruneau et al.'s 4R supply-chain-resilience
framework, evaluated by Mamdani fuzzy inference systems, solved by
NSGA-II, NRGA, and NSGA-III, and released as an open-source Python
reference implementation. The five contributions enumerated in
Section~\ref{sec:intro} are realised respectively in
Section~\ref{sec:resilience-mapping} (the 4R bijection),
Section~\ref{sec:tradeoff-analysis} (the H1 evidence that the 4-obj split
exposes a Robustness/Rapidity trade-off the 3-obj formulation cannot
represent),
Section~\ref{sec:performance} (the NSGA-III empirical comparison and the
32\% mean CPU advantage at unchanged Pareto-front quality),
Section~\ref{sec:sensitivity} (the rule-base and weight sensitivity
studies establishing $\le 1\%$ CV under realistic expert-elicitation
noise), and the open artifact at
\url{https://github.com/presidio-v/presidio-hardened-vol-assign}
(Zenodo DOI:~\url{https://doi.org/10.5281/zenodo.21083547}).

The aggregate finding is that the 4R framework is operationalisable
end-to-end in a humanitarian last-mile allocation model, that the
operationalisation produces empirical content beyond the theoretical
framing, and that the resulting recommendations are robust to the kinds
of expert-elicitation noise practitioners produce. For the
supply-chain-resilience research community, this closes a gap between
resilience theory and computational allocation practice. For the
civil-protection planner, this delivers a tool whose outputs read
directly in the vocabulary the planner already uses.

% =====================================================================
\section*{Data Availability Statement}
The full experimental data, including 540 H1+H2+H4 runs, 100 H3b samples,
and 66 H3a deletions, is reproducible from the open-source implementation
at \url{https://github.com/presidio-v/presidio-hardened-vol-assign}, tag
\texttt{v0.2.0}, archived on Zenodo (DOI:~\url{https://doi.org/10.5281/zenodo.21083547}).

% =====================================================================
\section*{Author Contributions}
% DRAFT CRediT — co-authors to confirm/adjust roles before submission.
Conceptualization, P.R., D.A.-A. and V.S.; methodology, P.R., D.A.-A. and V.S.;
software, P.R. and V.S.; validation, D.A.-A. and V.S.; formal analysis, V.S.;
investigation, P.R.; data curation, V.S.; writing---original draft preparation,
P.R. and V.S.; writing---review and editing, D.A.-A. and V.S.; visualization,
V.S.; supervision, D.A.-A. All authors have read and agreed to the published
version of the manuscript.

\section*{Funding}
This research received no external funding.

\section*{Institutional Review Board Statement}
Not applicable. This study used only synthetically generated problem instances
and involved no human participants or animals.

\section*{Informed Consent Statement}
Not applicable.

\section*{Acknowledgments}
The authors declare that no generative AI tools were used in the preparation of
this manuscript.

\section*{Conflicts of Interest}
The authors declare no conflicts of interest.

% =====================================================================
\appendix
\section{Full FIS rule bases}\label{app:rules}

This appendix enumerates every rule of the four fuzzy inference systems of the
four-objective formulation (FIS$_1$, FIS$_{2a}$, FIS$_{2b}$, FIS$_3$), followed
by the baseline FIS$_2$ used only as the H1 three-objective comparator. The
tables are generated directly from the rule-base constants in the open-source
implementation (\texttt{presidio\_vol\_assign.allocation.fis}); they are
authoritative and reproduce bit-exact under the released code.

\begin{table}[!ht]
\centering\footnotesize
\setlength{\tabcolsep}{4pt}
\caption{FIS$_1$ rule base: (VS, IDL, RTR) $\to$ ULPP, operationalising
Resourcefulness. Full $3^3 = 27$ Cartesian rule base.}
\label{tab:rules-fis1}
\begin{minipage}[t]{0.48\textwidth}\centering
\begin{tabular}{cccc}
\toprule
VS & IDL & RTR & ULPP \\
\midrule
Low & Minor & Short & Moderate \\
Low & Minor & Moderate & Low \\
Low & Minor & Long & Very Low \\
Low & Moderate & Short & High \\
Low & Moderate & Moderate & Moderate \\
Low & Moderate & Long & Low \\
Low & Severe & Short & Very High \\
Low & Severe & Moderate & High \\
Low & Severe & Long & Moderate \\
Medium & Minor & Short & High \\
Medium & Minor & Moderate & Moderate \\
Medium & Minor & Long & Low \\
Medium & Moderate & Short & Very High \\
Medium & Moderate & Moderate & High \\
\bottomrule
\end{tabular}
\end{minipage}\hfill
\begin{minipage}[t]{0.48\textwidth}\centering
\begin{tabular}{cccc}
\toprule
VS & IDL & RTR & ULPP \\
\midrule
Medium & Moderate & Long & Moderate \\
Medium & Severe & Short & Very High \\
Medium & Severe & Moderate & Very High \\
Medium & Severe & Long & High \\
High & Minor & Short & Very High \\
High & Minor & Moderate & High \\
High & Minor & Long & Moderate \\
High & Moderate & Short & Very High \\
High & Moderate & Moderate & Very High \\
High & Moderate & Long & High \\
High & Severe & Short & Very High \\
High & Severe & Moderate & Very High \\
High & Severe & Long & Very High \\
& & & \\
\bottomrule
\end{tabular}
\end{minipage}
\end{table}

\begin{table}[!ht]
\centering\footnotesize
\setlength{\tabcolsep}{6pt}
\caption{FIS$_{2a}$ rule base: (RCS, PHS) $\to$ TRD, operationalising
Robustness ($3 \times 3 = 9$ rules); FIS$_{2b}$ rule base:
(TD) $\to$ RPD, operationalising Rapidity (3 rules).}
\label{tab:rules-fis2}
\begin{minipage}[t]{0.55\textwidth}\centering
\textbf{FIS$_{2a}$ (Robustness)}\\[2pt]
\begin{tabular}{ccc}
\toprule
RCS & PHS & TRD \\
\midrule
Low & Low & Very Low \\
Low & Moderate & Low \\
Low & High & Moderate \\
Moderate & Low & Low \\
Moderate & Moderate & Moderate \\
Moderate & High & High \\
High & Low & Moderate \\
High & Moderate & High \\
High & High & Very High \\
\bottomrule
\end{tabular}
\end{minipage}\hfill
\begin{minipage}[t]{0.4\textwidth}\centering
\textbf{FIS$_{2b}$ (Rapidity)}\\[2pt]
\begin{tabular}{cc}
\toprule
TD & RPD \\
\midrule
Short & Low \\
Moderate & Moderate \\
Long & High \\
\bottomrule
\end{tabular}
\end{minipage}
\end{table}

\begin{table}[!ht]
\centering\footnotesize
\setlength{\tabcolsep}{4pt}
\caption{FIS$_3$ rule base: (COR, RDR, TD) $\to$ CAIL, operationalising
Redundancy. Full $3^3 = 27$ Cartesian rule base.}
\label{tab:rules-fis3}
\begin{minipage}[t]{0.48\textwidth}\centering
\begin{tabular}{cccc}
\toprule
COR & RDR & TD & CAIL \\
\midrule
Low & Slow & Short & Very Low \\
Low & Slow & Moderate & Low \\
Low & Slow & Long & Moderate \\
Low & Moderate & Short & Low \\
Low & Moderate & Moderate & Moderate \\
Low & Moderate & Long & High \\
Low & Rapid & Short & Moderate \\
Low & Rapid & Moderate & High \\
Low & Rapid & Long & High \\
Moderate & Slow & Short & Low \\
Moderate & Slow & Moderate & Moderate \\
Moderate & Slow & Long & High \\
Moderate & Moderate & Short & Moderate \\
Moderate & Moderate & Moderate & Moderate \\
\bottomrule
\end{tabular}
\end{minipage}\hfill
\begin{minipage}[t]{0.48\textwidth}\centering
\begin{tabular}{cccc}
\toprule
COR & RDR & TD & CAIL \\
\midrule
Moderate & Moderate & Long & High \\
Moderate & Rapid & Short & High \\
Moderate & Rapid & Moderate & High \\
Moderate & Rapid & Long & Very High \\
High & Slow & Short & Moderate \\
High & Slow & Moderate & High \\
High & Slow & Long & High \\
High & Moderate & Short & High \\
High & Moderate & Moderate & High \\
High & Moderate & Long & Very High \\
High & Rapid & Short & Very High \\
High & Rapid & Moderate & Very High \\
High & Rapid & Long & Very High \\
& & & \\
\bottomrule
\end{tabular}
\end{minipage}
\end{table}

\begin{table}[!ht]
\centering\footnotesize
\setlength{\tabcolsep}{6pt}
\caption{Baseline FIS$_2$ rule base: (TD, RWS) $\to$ TIL, the fused
transportation-infeasibility classifier used only as the three-objective
H1 comparator (Section~\ref{sec:tradeoff-analysis}). 9 rules.}
\label{tab:rules-fis2-baseline}
\begin{tabular}{ccc}
\toprule
TD & RWS & TIL \\
\midrule
Short & Low & Moderate \\
Short & Moderate & Low \\
Short & High & Very Low \\
Moderate & Low & High \\
Moderate & Moderate & Moderate \\
Moderate & High & Low \\
Long & Low & Very High \\
Long & Moderate & High \\
Long & High & Moderate \\
\bottomrule
\end{tabular}
\end{table}

\section{Statistical test outputs (extended)}\label{app:stats}

This appendix gives the complete pairwise Mann--Whitney $U$ two-sided test
outputs underlying the summary in Section~\ref{sec:perf-stats}. The
Mann--Whitney $U$ test is identical to the Wilcoxon rank-sum test, the
significance test most multi-objective optimisation studies report; we use the
$U$ formulation throughout for consistency with the normality-gated test
selection in Section~\ref{sec:perf-setup}. For each
problem size and each of the five quality metrics (NNS, HV, SM, CPU time, and
the MID diagnostic), we report the $p$-value of every algorithm pair (NSGA-II, NRGA,
NSGA-III), over 30 repetitions per cell. Tables~\ref{tab:appb-4obj}
and~\ref{tab:appb-3obj} cover the four-objective and the three-objective
baseline formulations respectively. The values are computed directly from
the released run manifest
(\texttt{experiments/results/h1\_h2\_h4/manifest.csv}) and reproduce under
the open-source code.

\begin{table}[!ht]
\centering\footnotesize
\caption{Pairwise Mann--Whitney $U$ two-sided $p$-values per
(size, metric) on the \textbf{four-objective} formulation, 30 reps per
cell. $^{*}$ marks $p<0.05$. ``---'' denotes a degenerate comparison
(both samples constant, e.g.\ NNS at the 100-individual ceiling).}
\label{tab:appb-4obj}
\begin{tabular}{l l ccccc}
\toprule
Size & Pair & NNS & HV & SM & CPU & MID$^{\dagger}$ \\
\midrule
\multirow{3}{*}{small} & NSGA-II vs NSGA-III & --- & $0.027^{*}$ & $<0.001^{*}$ & $<0.001^{*}$ & $<0.001^{*}$ \\
 & NSGA-II vs NRGA & $0.334$ & $<0.001^{*}$ & $<0.001^{*}$ & $<0.001^{*}$ & $<0.001^{*}$ \\
 & NSGA-III vs NRGA & $0.334$ & $0.006^{*}$ & $0.066$ & $<0.001^{*}$ & $<0.001^{*}$ \\
\midrule
\multirow{3}{*}{medium} & NSGA-II vs NSGA-III & --- & $0.001^{*}$ & $<0.001^{*}$ & $<0.001^{*}$ & $<0.001^{*}$ \\
 & NSGA-II vs NRGA & $0.334$ & $<0.001^{*}$ & $<0.001^{*}$ & $<0.001^{*}$ & $<0.001^{*}$ \\
 & NSGA-III vs NRGA & $0.334$ & $<0.001^{*}$ & $0.003^{*}$ & $<0.001^{*}$ & $0.112$ \\
\midrule
\multirow{3}{*}{large} & NSGA-II vs NSGA-III & --- & $<0.001^{*}$ & $<0.001^{*}$ & $<0.001^{*}$ & $<0.001^{*}$ \\
 & NSGA-II vs NRGA & $0.334$ & $<0.001^{*}$ & $<0.001^{*}$ & $<0.001^{*}$ & $<0.001^{*}$ \\
 & NSGA-III vs NRGA & $0.334$ & $<0.001^{*}$ & $0.652$ & $<0.001^{*}$ & $0.181$ \\
\bottomrule
\end{tabular}

\smallskip
{\footnotesize $^{\dagger}$\,MID is a diagnostic, not a ranking metric
(see Table~\ref{tab:stats-summary}); its column is placed last accordingly.}
\end{table}

\begin{table}[!ht]
\centering\footnotesize
\caption{Pairwise Mann--Whitney $U$ two-sided $p$-values per
(size, metric) on the \textbf{three-objective} baseline formulation,
30 reps per cell. $^{*}$ marks $p<0.05$.}
\label{tab:appb-3obj}
\begin{tabular}{l l ccccc}
\toprule
Size & Pair & NNS & HV & SM & CPU & MID$^{\dagger}$ \\
\midrule
\multirow{3}{*}{small} & NSGA-II vs NSGA-III & $0.334$ & $0.348$ & $<0.001^{*}$ & $<0.001^{*}$ & $0.044^{*}$ \\
 & NSGA-II vs NRGA & $0.021^{*}$ & $<0.001^{*}$ & $<0.001^{*}$ & $<0.001^{*}$ & $<0.001^{*}$ \\
 & NSGA-III vs NRGA & $0.006^{*}$ & $<0.001^{*}$ & $0.002^{*}$ & $<0.001^{*}$ & $0.158$ \\
\midrule
\multirow{3}{*}{medium} & NSGA-II vs NSGA-III & --- & $0.284$ & $<0.001^{*}$ & $<0.001^{*}$ & $0.041^{*}$ \\
 & NSGA-II vs NRGA & $0.082$ & $<0.001^{*}$ & $<0.001^{*}$ & $<0.001^{*}$ & $0.042^{*}$ \\
 & NSGA-III vs NRGA & $0.082$ & $<0.001^{*}$ & $<0.001^{*}$ & $<0.001^{*}$ & $0.784$ \\
\midrule
\multirow{3}{*}{large} & NSGA-II vs NSGA-III & --- & $0.018^{*}$ & $<0.001^{*}$ & $<0.001^{*}$ & $0.853$ \\
 & NSGA-II vs NRGA & $0.022^{*}$ & $<0.001^{*}$ & $<0.001^{*}$ & $<0.001^{*}$ & $<0.001^{*}$ \\
 & NSGA-III vs NRGA & $0.022^{*}$ & $<0.001^{*}$ & $0.007^{*}$ & $<0.001^{*}$ & $<0.001^{*}$ \\
\bottomrule
\end{tabular}

\smallskip
{\footnotesize $^{\dagger}$\,MID is a diagnostic, not a ranking metric
(see Table~\ref{tab:stats-summary}); its column is placed last accordingly.}
\end{table}

% =====================================================================
\bibliographystyle{plain}
\bibliography{lit}

\end{document}
